\documentclass[11pt,a4paper]{article}
% ============================================
% GIFT v3.4 — Shared Preamble
% ============================================
% Usage: % ============================================
% GIFT v3.4 — Shared Preamble
% ============================================
% Usage: % ============================================
% GIFT v3.4 — Shared Preamble
% ============================================
% Usage: \input{gift_preamble} after \documentclass

% ENCODING & FONTS
\usepackage[utf8]{inputenc}
\usepackage[T1]{fontenc}
\usepackage{lmodern}

% PAGE LAYOUT (Golden Ratio)
\usepackage[margin=1.618cm, top=2.618cm, bottom=2.618cm]{geometry}

% ESSENTIAL PACKAGES
\usepackage{float}
\usepackage{caption}
\usepackage{subcaption}
\usepackage{setspace}
\usepackage{fancyhdr}
\usepackage{xcolor}
\usepackage{hyperref}
\usepackage{csquotes}
\usepackage{amsmath}
\usepackage{amssymb}
\usepackage{amsthm}
\usepackage{mathtools}
\usepackage{booktabs}
\usepackage{longtable}
\usepackage{array}
\usepackage{multirow}
\usepackage{tikz}
\usepackage{graphicx}
\usepackage{listings}
\usepackage{enumitem}
% Unicode fallbacks (pandoc artifacts that survive markdown→LaTeX)
\DeclareUnicodeCharacter{00B0}{\ensuremath{^\circ}}
\DeclareUnicodeCharacter{221A}{\ensuremath{\sqrt{\,}}}
\DeclareUnicodeCharacter{2713}{\ensuremath{\checkmark}}
\DeclareUnicodeCharacter{2717}{\ensuremath{\times}}
\DeclareUnicodeCharacter{2070}{\ensuremath{^{0}}}
\DeclareUnicodeCharacter{00B9}{\ensuremath{^{1}}}
\DeclareUnicodeCharacter{00B2}{\ensuremath{^{2}}}
\DeclareUnicodeCharacter{00B3}{\ensuremath{^{3}}}
\DeclareUnicodeCharacter{2074}{\ensuremath{^{4}}}
\DeclareUnicodeCharacter{2075}{\ensuremath{^{5}}}
\DeclareUnicodeCharacter{2076}{\ensuremath{^{6}}}
\DeclareUnicodeCharacter{2077}{\ensuremath{^{7}}}
\DeclareUnicodeCharacter{2078}{\ensuremath{^{8}}}
\DeclareUnicodeCharacter{2079}{\ensuremath{^{9}}}
\DeclareUnicodeCharacter{2071}{\ensuremath{^{i}}}
\DeclareUnicodeCharacter{2080}{\ensuremath{_{0}}}
\DeclareUnicodeCharacter{2081}{\ensuremath{_{1}}}
\DeclareUnicodeCharacter{2082}{\ensuremath{_{2}}}
\DeclareUnicodeCharacter{2083}{\ensuremath{_{3}}}
\DeclareUnicodeCharacter{2084}{\ensuremath{_{4}}}
\DeclareUnicodeCharacter{2085}{\ensuremath{_{5}}}
\DeclareUnicodeCharacter{2086}{\ensuremath{_{6}}}
\DeclareUnicodeCharacter{2087}{\ensuremath{_{7}}}
\DeclareUnicodeCharacter{2088}{\ensuremath{_{8}}}
\DeclareUnicodeCharacter{2089}{\ensuremath{_{9}}}
\DeclareUnicodeCharacter{208A}{\ensuremath{_{+}}}
\DeclareUnicodeCharacter{208B}{\ensuremath{_{-}}}
\DeclareUnicodeCharacter{2192}{\ensuremath{\to}}
\DeclareUnicodeCharacter{2502}{|}
\DeclareUnicodeCharacter{25BC}{\ensuremath{\blacktriangledown}}
\DeclareUnicodeCharacter{2500}{\ensuremath{-}}
\DeclareUnicodeCharacter{2514}{\ensuremath{\llcorner}}
\DeclareUnicodeCharacter{251C}{\ensuremath{\vdash}}
\DeclareUnicodeCharacter{2026}{\ldots}
\DeclareUnicodeCharacter{2208}{\ensuremath{\in}}
\DeclareUnicodeCharacter{2229}{\ensuremath{\cap}}
\DeclareUnicodeCharacter{222A}{\ensuremath{\cup}}
\DeclareUnicodeCharacter{2295}{\ensuremath{\oplus}}
\DeclareUnicodeCharacter{2297}{\ensuremath{\otimes}}
\DeclareUnicodeCharacter{2245}{\ensuremath{\cong}}
\DeclareUnicodeCharacter{2248}{\ensuremath{\approx}}
\DeclareUnicodeCharacter{2260}{\ensuremath{\neq}}
\DeclareUnicodeCharacter{2261}{\ensuremath{\equiv}}
\DeclareUnicodeCharacter{2264}{\ensuremath{\leq}}
\DeclareUnicodeCharacter{2265}{\ensuremath{\geq}}
\DeclareUnicodeCharacter{00D7}{\ensuremath{\times}}
\DeclareUnicodeCharacter{00B1}{\ensuremath{\pm}}
\DeclareUnicodeCharacter{2225}{\ensuremath{\|}}
% Greek letters used in body text (pandoc passes them through unchanged)
\DeclareUnicodeCharacter{03B1}{\ensuremath{\alpha}}
\DeclareUnicodeCharacter{03B2}{\ensuremath{\beta}}
\DeclareUnicodeCharacter{03B3}{\ensuremath{\gamma}}
\DeclareUnicodeCharacter{03B4}{\ensuremath{\delta}}
\DeclareUnicodeCharacter{03B5}{\ensuremath{\varepsilon}}
\DeclareUnicodeCharacter{03B6}{\ensuremath{\zeta}}
\DeclareUnicodeCharacter{03B7}{\ensuremath{\eta}}
\DeclareUnicodeCharacter{03B8}{\ensuremath{\theta}}
\DeclareUnicodeCharacter{03B9}{\ensuremath{\iota}}
\DeclareUnicodeCharacter{03BA}{\ensuremath{\kappa}}
\DeclareUnicodeCharacter{03BB}{\ensuremath{\lambda}}
\DeclareUnicodeCharacter{03BC}{\ensuremath{\mu}}
\DeclareUnicodeCharacter{03BD}{\ensuremath{\nu}}
\DeclareUnicodeCharacter{03BE}{\ensuremath{\xi}}
\DeclareUnicodeCharacter{03C0}{\ensuremath{\pi}}
\DeclareUnicodeCharacter{03C1}{\ensuremath{\rho}}
\DeclareUnicodeCharacter{03C3}{\ensuremath{\sigma}}
\DeclareUnicodeCharacter{03C4}{\ensuremath{\tau}}
\DeclareUnicodeCharacter{03C5}{\ensuremath{\upsilon}}
\DeclareUnicodeCharacter{03C6}{\ensuremath{\varphi}}
\DeclareUnicodeCharacter{03C7}{\ensuremath{\chi}}
\DeclareUnicodeCharacter{03C8}{\ensuremath{\psi}}
\DeclareUnicodeCharacter{03C9}{\ensuremath{\omega}}
\DeclareUnicodeCharacter{0394}{\ensuremath{\Delta}}
\DeclareUnicodeCharacter{0398}{\ensuremath{\Theta}}
\DeclareUnicodeCharacter{039B}{\ensuremath{\Lambda}}
\DeclareUnicodeCharacter{03A0}{\ensuremath{\Pi}}
\DeclareUnicodeCharacter{03A3}{\ensuremath{\Sigma}}
\DeclareUnicodeCharacter{03A6}{\ensuremath{\Phi}}
\DeclareUnicodeCharacter{03A8}{\ensuremath{\Psi}}
\DeclareUnicodeCharacter{03A9}{\ensuremath{\Omega}}
\DeclareUnicodeCharacter{2200}{\ensuremath{\forall}}
\DeclareUnicodeCharacter{2203}{\ensuremath{\exists}}
\DeclareUnicodeCharacter{2207}{\ensuremath{\nabla}}
\DeclareUnicodeCharacter{2202}{\ensuremath{\partial}}
\DeclareUnicodeCharacter{2218}{\ensuremath{\circ}}
\DeclareUnicodeCharacter{2032}{\ensuremath{\prime}}
\DeclareUnicodeCharacter{2113}{\ensuremath{\ell}}
\DeclareUnicodeCharacter{210F}{\ensuremath{\hbar}}
\DeclareUnicodeCharacter{27E8}{\ensuremath{\langle}}
\DeclareUnicodeCharacter{27E9}{\ensuremath{\rangle}}
\DeclareUnicodeCharacter{2308}{\ensuremath{\lceil}}
\DeclareUnicodeCharacter{2309}{\ensuremath{\rceil}}
\DeclareUnicodeCharacter{230A}{\ensuremath{\lfloor}}
\DeclareUnicodeCharacter{230B}{\ensuremath{\rfloor}}
\DeclareUnicodeCharacter{222B}{\ensuremath{\int}}
\DeclareUnicodeCharacter{2211}{\ensuremath{\sum}}
\DeclareUnicodeCharacter{220F}{\ensuremath{\prod}}
\DeclareUnicodeCharacter{221E}{\ensuremath{\infty}}
\DeclareUnicodeCharacter{00B5}{\ensuremath{\mu}}
\DeclareUnicodeCharacter{2236}{\ensuremath{\colon}}
\DeclareUnicodeCharacter{2192}{\ensuremath{\to}}
\DeclareUnicodeCharacter{21A6}{\ensuremath{\mapsto}}
\DeclareUnicodeCharacter{21D2}{\ensuremath{\Rightarrow}}
\DeclareUnicodeCharacter{21D4}{\ensuremath{\Leftrightarrow}}
\DeclareUnicodeCharacter{2282}{\ensuremath{\subset}}
\DeclareUnicodeCharacter{2286}{\ensuremath{\subseteq}}
\DeclareUnicodeCharacter{2287}{\ensuremath{\supseteq}}
\DeclareUnicodeCharacter{222B}{\ensuremath{\int}}
\DeclareUnicodeCharacter{2205}{\ensuremath{\emptyset}}
\DeclareUnicodeCharacter{2299}{\ensuremath{\odot}}
\DeclareUnicodeCharacter{2220}{\ensuremath{\angle}}
\DeclareUnicodeCharacter{2135}{\ensuremath{\aleph}}
\DeclareUnicodeCharacter{207B}{\ensuremath{^{-}}}
\DeclareUnicodeCharacter{207A}{\ensuremath{^{+}}}
\DeclareUnicodeCharacter{220E}{\ensuremath{\blacksquare}}
\DeclareUnicodeCharacter{2061}{}
\DeclareUnicodeCharacter{2062}{}
\DeclareUnicodeCharacter{2063}{}
\DeclareUnicodeCharacter{2009}{\,}
\DeclareUnicodeCharacter{200A}{\,}
\DeclareUnicodeCharacter{200B}{}
\DeclareUnicodeCharacter{2212}{\ensuremath{-}}
\DeclareUnicodeCharacter{2194}{\ensuremath{\leftrightarrow}}
\DeclareUnicodeCharacter{2099}{\ensuremath{_{n}}}
\DeclareUnicodeCharacter{2098}{\ensuremath{_{m}}}
\DeclareUnicodeCharacter{2096}{\ensuremath{_{k}}}
\DeclareUnicodeCharacter{2095}{\ensuremath{_{h}}}
\DeclareUnicodeCharacter{2090}{\ensuremath{_{a}}}
\DeclareUnicodeCharacter{2091}{\ensuremath{_{e}}}
\DeclareUnicodeCharacter{2093}{\ensuremath{_{x}}}
\DeclareUnicodeCharacter{1D62}{\ensuremath{_{i}}}
\DeclareUnicodeCharacter{2C7C}{\ensuremath{_{j}}}
\DeclareUnicodeCharacter{2C7B}{\ensuremath{^{e}}}
\DeclareUnicodeCharacter{1D52}{\ensuremath{^{o}}}
\DeclareUnicodeCharacter{1D5}{\ensuremath{^{u}}}
\DeclareUnicodeCharacter{02E1}{\ensuremath{^{l}}}
\DeclareUnicodeCharacter{2295}{\ensuremath{\oplus}}
\DeclareUnicodeCharacter{27FA}{\ensuremath{\Longleftrightarrow}}
\DeclareUnicodeCharacter{27F9}{\ensuremath{\Longrightarrow}}
\DeclareUnicodeCharacter{2299}{\ensuremath{\odot}}
% Provide \tightlist used by pandoc-generated lists (no-op)
\providecommand{\tightlist}{\setlength{\itemsep}{0pt}\setlength{\parskip}{0pt}}
% pandoc longtable column specs: provide \real as passthrough and a
% `none` counter so that calc's \real{x} expansion path works.
\newcounter{none}
\providecommand{\real}[1]{#1}

% LISTINGS
\lstset{
    basicstyle=\small\ttfamily,
    breaklines=true, frame=single,
    keepspaces=true, showstringspaces=false,
    aboveskip=0.8em, belowskip=0.8em
}

% HEADER/FOOTER
\setlength{\headheight}{14pt}
\pagestyle{fancy}
\fancyhf{}
\fancyhead[R]{\thepage}
\renewcommand{\headrulewidth}{0.2pt}

% HYPERREF
\hypersetup{
    colorlinks=true,
    linkcolor=blue, citecolor=blue, urlcolor=blue,
    pdfauthor={Brieuc de La Fourniere}
}

% SPACING
\setstretch{1.2}
\setlength{\parskip}{0.4em}
\setlength{\parindent}{0pt}

% TITLE FORMATTING
\usepackage{titling}
\pretitle{\LARGE\bfseries}
\posttitle{\vspace{-0.4em}}
\preauthor{}
\postauthor{}
\predate{}
\postdate{}
\setlength{\droptitle}{-2.0em}

% CUSTOM COMMANDS — GIFT notation
\newcommand{\E}{\mathrm{E}}
\newcommand{\Gtwo}{\mathrm{G}_2}
\newcommand{\Kseven}{K_7}
\newcommand{\AdS}{\mathrm{AdS}}
\newcommand{\SU}{\mathrm{SU}}
\newcommand{\SO}{\mathrm{SO}}
\newcommand{\U}{\mathrm{U}}
\newcommand{\Sp}{\mathrm{Sp}}
\newcommand{\PSL}{\mathrm{PSL}}
\newcommand{\SM}{\mathrm{SM}}
\newcommand{\Tr}{\mathrm{Tr}}
\newcommand{\rank}{\mathrm{rank}}
\newcommand{\Vol}{\mathrm{Vol}}
\newcommand{\Hol}{\mathrm{Hol}}
\newcommand{\Ric}{\mathrm{Ric}}
\newcommand{\Rm}{\mathrm{Rm}}
\newcommand{\Scal}{\mathrm{Scal}}
\DeclareMathOperator{\diag}{diag}
\DeclareMathOperator{\cond}{cond}
\newcommand{\GIFT}{\textrm{GIFT}}
\newcommand{\CP}{\mathrm{CP}}
\newcommand{\CKM}{\mathrm{CKM}}
\newcommand{\PMNS}{\mathrm{PMNS}}
\newcommand{\EW}{\mathrm{EW}}
\newcommand{\Pl}{\mathrm{Pl}}
\newcommand{\DE}{\mathrm{DE}}
\newcommand{\DM}{\mathrm{DM}}
\newcommand{\KK}{\mathrm{KK}}
\newcommand{\NK}{\mathrm{NK}}
\newcommand{\TCS}{\mathrm{TCS}}
\newcommand{\PINN}{\mathrm{PINN}}
\newcommand{\btwo}{b_2}
\newcommand{\bthree}{b_3}
\newcommand{\typeI}{\textsc{I}}
\newcommand{\typeII}{\textsc{II}}
\newcommand{\typeIII}{\textsc{III}}
\newcommand{\typeIV}{\textsc{IV}}
\newcommand{\proven}{\textsc{Proven}}

% PDF bookmark safety
\pdfstringdefDisableCommands{%
  \def\textsubscript#1{#1}%
  \def\textsuperscript#1{#1}%
  \def\CP{CP}%
  \def\Gtwo{G2}%
  \def\Kseven{K7}%
  \def\GIFT{GIFT}%
  \def\E{E}%
  \def\SM{SM}%
}

% THEOREM ENVIRONMENTS
\newtheorem{theorem}{Theorem}[section]
\newtheorem{proposition}[theorem]{Proposition}
\newtheorem{lemma}[theorem]{Lemma}
\newtheorem{corollary}[theorem]{Corollary}
\theoremstyle{definition}
\newtheorem{definition}[theorem]{Definition}
\theoremstyle{remark}
\newtheorem{remark}[theorem]{Remark}
 after \documentclass

% ENCODING & FONTS
\usepackage[utf8]{inputenc}
\usepackage[T1]{fontenc}
\usepackage{lmodern}

% PAGE LAYOUT (Golden Ratio)
\usepackage[margin=1.618cm, top=2.618cm, bottom=2.618cm]{geometry}

% ESSENTIAL PACKAGES
\usepackage{float}
\usepackage{caption}
\usepackage{subcaption}
\usepackage{setspace}
\usepackage{fancyhdr}
\usepackage{xcolor}
\usepackage{hyperref}
\usepackage{csquotes}
\usepackage{amsmath}
\usepackage{amssymb}
\usepackage{amsthm}
\usepackage{mathtools}
\usepackage{booktabs}
\usepackage{longtable}
\usepackage{array}
\usepackage{multirow}
\usepackage{tikz}
\usepackage{graphicx}
\usepackage{listings}
\usepackage{enumitem}
% Unicode fallbacks (pandoc artifacts that survive markdown→LaTeX)
\DeclareUnicodeCharacter{00B0}{\ensuremath{^\circ}}
\DeclareUnicodeCharacter{221A}{\ensuremath{\sqrt{\,}}}
\DeclareUnicodeCharacter{2713}{\ensuremath{\checkmark}}
\DeclareUnicodeCharacter{2717}{\ensuremath{\times}}
\DeclareUnicodeCharacter{2070}{\ensuremath{^{0}}}
\DeclareUnicodeCharacter{00B9}{\ensuremath{^{1}}}
\DeclareUnicodeCharacter{00B2}{\ensuremath{^{2}}}
\DeclareUnicodeCharacter{00B3}{\ensuremath{^{3}}}
\DeclareUnicodeCharacter{2074}{\ensuremath{^{4}}}
\DeclareUnicodeCharacter{2075}{\ensuremath{^{5}}}
\DeclareUnicodeCharacter{2076}{\ensuremath{^{6}}}
\DeclareUnicodeCharacter{2077}{\ensuremath{^{7}}}
\DeclareUnicodeCharacter{2078}{\ensuremath{^{8}}}
\DeclareUnicodeCharacter{2079}{\ensuremath{^{9}}}
\DeclareUnicodeCharacter{2071}{\ensuremath{^{i}}}
\DeclareUnicodeCharacter{2080}{\ensuremath{_{0}}}
\DeclareUnicodeCharacter{2081}{\ensuremath{_{1}}}
\DeclareUnicodeCharacter{2082}{\ensuremath{_{2}}}
\DeclareUnicodeCharacter{2083}{\ensuremath{_{3}}}
\DeclareUnicodeCharacter{2084}{\ensuremath{_{4}}}
\DeclareUnicodeCharacter{2085}{\ensuremath{_{5}}}
\DeclareUnicodeCharacter{2086}{\ensuremath{_{6}}}
\DeclareUnicodeCharacter{2087}{\ensuremath{_{7}}}
\DeclareUnicodeCharacter{2088}{\ensuremath{_{8}}}
\DeclareUnicodeCharacter{2089}{\ensuremath{_{9}}}
\DeclareUnicodeCharacter{208A}{\ensuremath{_{+}}}
\DeclareUnicodeCharacter{208B}{\ensuremath{_{-}}}
\DeclareUnicodeCharacter{2192}{\ensuremath{\to}}
\DeclareUnicodeCharacter{2502}{|}
\DeclareUnicodeCharacter{25BC}{\ensuremath{\blacktriangledown}}
\DeclareUnicodeCharacter{2500}{\ensuremath{-}}
\DeclareUnicodeCharacter{2514}{\ensuremath{\llcorner}}
\DeclareUnicodeCharacter{251C}{\ensuremath{\vdash}}
\DeclareUnicodeCharacter{2026}{\ldots}
\DeclareUnicodeCharacter{2208}{\ensuremath{\in}}
\DeclareUnicodeCharacter{2229}{\ensuremath{\cap}}
\DeclareUnicodeCharacter{222A}{\ensuremath{\cup}}
\DeclareUnicodeCharacter{2295}{\ensuremath{\oplus}}
\DeclareUnicodeCharacter{2297}{\ensuremath{\otimes}}
\DeclareUnicodeCharacter{2245}{\ensuremath{\cong}}
\DeclareUnicodeCharacter{2248}{\ensuremath{\approx}}
\DeclareUnicodeCharacter{2260}{\ensuremath{\neq}}
\DeclareUnicodeCharacter{2261}{\ensuremath{\equiv}}
\DeclareUnicodeCharacter{2264}{\ensuremath{\leq}}
\DeclareUnicodeCharacter{2265}{\ensuremath{\geq}}
\DeclareUnicodeCharacter{00D7}{\ensuremath{\times}}
\DeclareUnicodeCharacter{00B1}{\ensuremath{\pm}}
\DeclareUnicodeCharacter{2225}{\ensuremath{\|}}
% Greek letters used in body text (pandoc passes them through unchanged)
\DeclareUnicodeCharacter{03B1}{\ensuremath{\alpha}}
\DeclareUnicodeCharacter{03B2}{\ensuremath{\beta}}
\DeclareUnicodeCharacter{03B3}{\ensuremath{\gamma}}
\DeclareUnicodeCharacter{03B4}{\ensuremath{\delta}}
\DeclareUnicodeCharacter{03B5}{\ensuremath{\varepsilon}}
\DeclareUnicodeCharacter{03B6}{\ensuremath{\zeta}}
\DeclareUnicodeCharacter{03B7}{\ensuremath{\eta}}
\DeclareUnicodeCharacter{03B8}{\ensuremath{\theta}}
\DeclareUnicodeCharacter{03B9}{\ensuremath{\iota}}
\DeclareUnicodeCharacter{03BA}{\ensuremath{\kappa}}
\DeclareUnicodeCharacter{03BB}{\ensuremath{\lambda}}
\DeclareUnicodeCharacter{03BC}{\ensuremath{\mu}}
\DeclareUnicodeCharacter{03BD}{\ensuremath{\nu}}
\DeclareUnicodeCharacter{03BE}{\ensuremath{\xi}}
\DeclareUnicodeCharacter{03C0}{\ensuremath{\pi}}
\DeclareUnicodeCharacter{03C1}{\ensuremath{\rho}}
\DeclareUnicodeCharacter{03C3}{\ensuremath{\sigma}}
\DeclareUnicodeCharacter{03C4}{\ensuremath{\tau}}
\DeclareUnicodeCharacter{03C5}{\ensuremath{\upsilon}}
\DeclareUnicodeCharacter{03C6}{\ensuremath{\varphi}}
\DeclareUnicodeCharacter{03C7}{\ensuremath{\chi}}
\DeclareUnicodeCharacter{03C8}{\ensuremath{\psi}}
\DeclareUnicodeCharacter{03C9}{\ensuremath{\omega}}
\DeclareUnicodeCharacter{0394}{\ensuremath{\Delta}}
\DeclareUnicodeCharacter{0398}{\ensuremath{\Theta}}
\DeclareUnicodeCharacter{039B}{\ensuremath{\Lambda}}
\DeclareUnicodeCharacter{03A0}{\ensuremath{\Pi}}
\DeclareUnicodeCharacter{03A3}{\ensuremath{\Sigma}}
\DeclareUnicodeCharacter{03A6}{\ensuremath{\Phi}}
\DeclareUnicodeCharacter{03A8}{\ensuremath{\Psi}}
\DeclareUnicodeCharacter{03A9}{\ensuremath{\Omega}}
\DeclareUnicodeCharacter{2200}{\ensuremath{\forall}}
\DeclareUnicodeCharacter{2203}{\ensuremath{\exists}}
\DeclareUnicodeCharacter{2207}{\ensuremath{\nabla}}
\DeclareUnicodeCharacter{2202}{\ensuremath{\partial}}
\DeclareUnicodeCharacter{2218}{\ensuremath{\circ}}
\DeclareUnicodeCharacter{2032}{\ensuremath{\prime}}
\DeclareUnicodeCharacter{2113}{\ensuremath{\ell}}
\DeclareUnicodeCharacter{210F}{\ensuremath{\hbar}}
\DeclareUnicodeCharacter{27E8}{\ensuremath{\langle}}
\DeclareUnicodeCharacter{27E9}{\ensuremath{\rangle}}
\DeclareUnicodeCharacter{2308}{\ensuremath{\lceil}}
\DeclareUnicodeCharacter{2309}{\ensuremath{\rceil}}
\DeclareUnicodeCharacter{230A}{\ensuremath{\lfloor}}
\DeclareUnicodeCharacter{230B}{\ensuremath{\rfloor}}
\DeclareUnicodeCharacter{222B}{\ensuremath{\int}}
\DeclareUnicodeCharacter{2211}{\ensuremath{\sum}}
\DeclareUnicodeCharacter{220F}{\ensuremath{\prod}}
\DeclareUnicodeCharacter{221E}{\ensuremath{\infty}}
\DeclareUnicodeCharacter{00B5}{\ensuremath{\mu}}
\DeclareUnicodeCharacter{2236}{\ensuremath{\colon}}
\DeclareUnicodeCharacter{2192}{\ensuremath{\to}}
\DeclareUnicodeCharacter{21A6}{\ensuremath{\mapsto}}
\DeclareUnicodeCharacter{21D2}{\ensuremath{\Rightarrow}}
\DeclareUnicodeCharacter{21D4}{\ensuremath{\Leftrightarrow}}
\DeclareUnicodeCharacter{2282}{\ensuremath{\subset}}
\DeclareUnicodeCharacter{2286}{\ensuremath{\subseteq}}
\DeclareUnicodeCharacter{2287}{\ensuremath{\supseteq}}
\DeclareUnicodeCharacter{222B}{\ensuremath{\int}}
\DeclareUnicodeCharacter{2205}{\ensuremath{\emptyset}}
\DeclareUnicodeCharacter{2299}{\ensuremath{\odot}}
\DeclareUnicodeCharacter{2220}{\ensuremath{\angle}}
\DeclareUnicodeCharacter{2135}{\ensuremath{\aleph}}
\DeclareUnicodeCharacter{207B}{\ensuremath{^{-}}}
\DeclareUnicodeCharacter{207A}{\ensuremath{^{+}}}
\DeclareUnicodeCharacter{220E}{\ensuremath{\blacksquare}}
\DeclareUnicodeCharacter{2061}{}
\DeclareUnicodeCharacter{2062}{}
\DeclareUnicodeCharacter{2063}{}
\DeclareUnicodeCharacter{2009}{\,}
\DeclareUnicodeCharacter{200A}{\,}
\DeclareUnicodeCharacter{200B}{}
\DeclareUnicodeCharacter{2212}{\ensuremath{-}}
\DeclareUnicodeCharacter{2194}{\ensuremath{\leftrightarrow}}
\DeclareUnicodeCharacter{2099}{\ensuremath{_{n}}}
\DeclareUnicodeCharacter{2098}{\ensuremath{_{m}}}
\DeclareUnicodeCharacter{2096}{\ensuremath{_{k}}}
\DeclareUnicodeCharacter{2095}{\ensuremath{_{h}}}
\DeclareUnicodeCharacter{2090}{\ensuremath{_{a}}}
\DeclareUnicodeCharacter{2091}{\ensuremath{_{e}}}
\DeclareUnicodeCharacter{2093}{\ensuremath{_{x}}}
\DeclareUnicodeCharacter{1D62}{\ensuremath{_{i}}}
\DeclareUnicodeCharacter{2C7C}{\ensuremath{_{j}}}
\DeclareUnicodeCharacter{2C7B}{\ensuremath{^{e}}}
\DeclareUnicodeCharacter{1D52}{\ensuremath{^{o}}}
\DeclareUnicodeCharacter{1D5}{\ensuremath{^{u}}}
\DeclareUnicodeCharacter{02E1}{\ensuremath{^{l}}}
\DeclareUnicodeCharacter{2295}{\ensuremath{\oplus}}
\DeclareUnicodeCharacter{27FA}{\ensuremath{\Longleftrightarrow}}
\DeclareUnicodeCharacter{27F9}{\ensuremath{\Longrightarrow}}
\DeclareUnicodeCharacter{2299}{\ensuremath{\odot}}
% Provide \tightlist used by pandoc-generated lists (no-op)
\providecommand{\tightlist}{\setlength{\itemsep}{0pt}\setlength{\parskip}{0pt}}
% pandoc longtable column specs: provide \real as passthrough and a
% `none` counter so that calc's \real{x} expansion path works.
\newcounter{none}
\providecommand{\real}[1]{#1}

% LISTINGS
\lstset{
    basicstyle=\small\ttfamily,
    breaklines=true, frame=single,
    keepspaces=true, showstringspaces=false,
    aboveskip=0.8em, belowskip=0.8em
}

% HEADER/FOOTER
\setlength{\headheight}{14pt}
\pagestyle{fancy}
\fancyhf{}
\fancyhead[R]{\thepage}
\renewcommand{\headrulewidth}{0.2pt}

% HYPERREF
\hypersetup{
    colorlinks=true,
    linkcolor=blue, citecolor=blue, urlcolor=blue,
    pdfauthor={Brieuc de La Fourniere}
}

% SPACING
\setstretch{1.2}
\setlength{\parskip}{0.4em}
\setlength{\parindent}{0pt}

% TITLE FORMATTING
\usepackage{titling}
\pretitle{\LARGE\bfseries}
\posttitle{\vspace{-0.4em}}
\preauthor{}
\postauthor{}
\predate{}
\postdate{}
\setlength{\droptitle}{-2.0em}

% CUSTOM COMMANDS — GIFT notation
\newcommand{\E}{\mathrm{E}}
\newcommand{\Gtwo}{\mathrm{G}_2}
\newcommand{\Kseven}{K_7}
\newcommand{\AdS}{\mathrm{AdS}}
\newcommand{\SU}{\mathrm{SU}}
\newcommand{\SO}{\mathrm{SO}}
\newcommand{\U}{\mathrm{U}}
\newcommand{\Sp}{\mathrm{Sp}}
\newcommand{\PSL}{\mathrm{PSL}}
\newcommand{\SM}{\mathrm{SM}}
\newcommand{\Tr}{\mathrm{Tr}}
\newcommand{\rank}{\mathrm{rank}}
\newcommand{\Vol}{\mathrm{Vol}}
\newcommand{\Hol}{\mathrm{Hol}}
\newcommand{\Ric}{\mathrm{Ric}}
\newcommand{\Rm}{\mathrm{Rm}}
\newcommand{\Scal}{\mathrm{Scal}}
\DeclareMathOperator{\diag}{diag}
\DeclareMathOperator{\cond}{cond}
\newcommand{\GIFT}{\textrm{GIFT}}
\newcommand{\CP}{\mathrm{CP}}
\newcommand{\CKM}{\mathrm{CKM}}
\newcommand{\PMNS}{\mathrm{PMNS}}
\newcommand{\EW}{\mathrm{EW}}
\newcommand{\Pl}{\mathrm{Pl}}
\newcommand{\DE}{\mathrm{DE}}
\newcommand{\DM}{\mathrm{DM}}
\newcommand{\KK}{\mathrm{KK}}
\newcommand{\NK}{\mathrm{NK}}
\newcommand{\TCS}{\mathrm{TCS}}
\newcommand{\PINN}{\mathrm{PINN}}
\newcommand{\btwo}{b_2}
\newcommand{\bthree}{b_3}
\newcommand{\typeI}{\textsc{I}}
\newcommand{\typeII}{\textsc{II}}
\newcommand{\typeIII}{\textsc{III}}
\newcommand{\typeIV}{\textsc{IV}}
\newcommand{\proven}{\textsc{Proven}}

% PDF bookmark safety
\pdfstringdefDisableCommands{%
  \def\textsubscript#1{#1}%
  \def\textsuperscript#1{#1}%
  \def\CP{CP}%
  \def\Gtwo{G2}%
  \def\Kseven{K7}%
  \def\GIFT{GIFT}%
  \def\E{E}%
  \def\SM{SM}%
}

% THEOREM ENVIRONMENTS
\newtheorem{theorem}{Theorem}[section]
\newtheorem{proposition}[theorem]{Proposition}
\newtheorem{lemma}[theorem]{Lemma}
\newtheorem{corollary}[theorem]{Corollary}
\theoremstyle{definition}
\newtheorem{definition}[theorem]{Definition}
\theoremstyle{remark}
\newtheorem{remark}[theorem]{Remark}
 after \documentclass

% ENCODING & FONTS
\usepackage[utf8]{inputenc}
\usepackage[T1]{fontenc}
\usepackage{lmodern}

% PAGE LAYOUT (Golden Ratio)
\usepackage[margin=1.618cm, top=2.618cm, bottom=2.618cm]{geometry}

% ESSENTIAL PACKAGES
\usepackage{float}
\usepackage{caption}
\usepackage{subcaption}
\usepackage{setspace}
\usepackage{fancyhdr}
\usepackage{xcolor}
\usepackage{hyperref}
\usepackage{csquotes}
\usepackage{amsmath}
\usepackage{amssymb}
\usepackage{amsthm}
\usepackage{mathtools}
\usepackage{booktabs}
\usepackage{longtable}
\usepackage{array}
\usepackage{multirow}
\usepackage{tikz}
\usepackage{graphicx}
\usepackage{listings}
\usepackage{enumitem}
% Unicode fallbacks (pandoc artifacts that survive markdown→LaTeX)
\DeclareUnicodeCharacter{00B0}{\ensuremath{^\circ}}
\DeclareUnicodeCharacter{221A}{\ensuremath{\sqrt{\,}}}
\DeclareUnicodeCharacter{2713}{\ensuremath{\checkmark}}
\DeclareUnicodeCharacter{2717}{\ensuremath{\times}}
\DeclareUnicodeCharacter{2070}{\ensuremath{^{0}}}
\DeclareUnicodeCharacter{00B9}{\ensuremath{^{1}}}
\DeclareUnicodeCharacter{00B2}{\ensuremath{^{2}}}
\DeclareUnicodeCharacter{00B3}{\ensuremath{^{3}}}
\DeclareUnicodeCharacter{2074}{\ensuremath{^{4}}}
\DeclareUnicodeCharacter{2075}{\ensuremath{^{5}}}
\DeclareUnicodeCharacter{2076}{\ensuremath{^{6}}}
\DeclareUnicodeCharacter{2077}{\ensuremath{^{7}}}
\DeclareUnicodeCharacter{2078}{\ensuremath{^{8}}}
\DeclareUnicodeCharacter{2079}{\ensuremath{^{9}}}
\DeclareUnicodeCharacter{2071}{\ensuremath{^{i}}}
\DeclareUnicodeCharacter{2080}{\ensuremath{_{0}}}
\DeclareUnicodeCharacter{2081}{\ensuremath{_{1}}}
\DeclareUnicodeCharacter{2082}{\ensuremath{_{2}}}
\DeclareUnicodeCharacter{2083}{\ensuremath{_{3}}}
\DeclareUnicodeCharacter{2084}{\ensuremath{_{4}}}
\DeclareUnicodeCharacter{2085}{\ensuremath{_{5}}}
\DeclareUnicodeCharacter{2086}{\ensuremath{_{6}}}
\DeclareUnicodeCharacter{2087}{\ensuremath{_{7}}}
\DeclareUnicodeCharacter{2088}{\ensuremath{_{8}}}
\DeclareUnicodeCharacter{2089}{\ensuremath{_{9}}}
\DeclareUnicodeCharacter{208A}{\ensuremath{_{+}}}
\DeclareUnicodeCharacter{208B}{\ensuremath{_{-}}}
\DeclareUnicodeCharacter{2192}{\ensuremath{\to}}
\DeclareUnicodeCharacter{2502}{|}
\DeclareUnicodeCharacter{25BC}{\ensuremath{\blacktriangledown}}
\DeclareUnicodeCharacter{2500}{\ensuremath{-}}
\DeclareUnicodeCharacter{2514}{\ensuremath{\llcorner}}
\DeclareUnicodeCharacter{251C}{\ensuremath{\vdash}}
\DeclareUnicodeCharacter{2026}{\ldots}
\DeclareUnicodeCharacter{2208}{\ensuremath{\in}}
\DeclareUnicodeCharacter{2229}{\ensuremath{\cap}}
\DeclareUnicodeCharacter{222A}{\ensuremath{\cup}}
\DeclareUnicodeCharacter{2295}{\ensuremath{\oplus}}
\DeclareUnicodeCharacter{2297}{\ensuremath{\otimes}}
\DeclareUnicodeCharacter{2245}{\ensuremath{\cong}}
\DeclareUnicodeCharacter{2248}{\ensuremath{\approx}}
\DeclareUnicodeCharacter{2260}{\ensuremath{\neq}}
\DeclareUnicodeCharacter{2261}{\ensuremath{\equiv}}
\DeclareUnicodeCharacter{2264}{\ensuremath{\leq}}
\DeclareUnicodeCharacter{2265}{\ensuremath{\geq}}
\DeclareUnicodeCharacter{00D7}{\ensuremath{\times}}
\DeclareUnicodeCharacter{00B1}{\ensuremath{\pm}}
\DeclareUnicodeCharacter{2225}{\ensuremath{\|}}
% Greek letters used in body text (pandoc passes them through unchanged)
\DeclareUnicodeCharacter{03B1}{\ensuremath{\alpha}}
\DeclareUnicodeCharacter{03B2}{\ensuremath{\beta}}
\DeclareUnicodeCharacter{03B3}{\ensuremath{\gamma}}
\DeclareUnicodeCharacter{03B4}{\ensuremath{\delta}}
\DeclareUnicodeCharacter{03B5}{\ensuremath{\varepsilon}}
\DeclareUnicodeCharacter{03B6}{\ensuremath{\zeta}}
\DeclareUnicodeCharacter{03B7}{\ensuremath{\eta}}
\DeclareUnicodeCharacter{03B8}{\ensuremath{\theta}}
\DeclareUnicodeCharacter{03B9}{\ensuremath{\iota}}
\DeclareUnicodeCharacter{03BA}{\ensuremath{\kappa}}
\DeclareUnicodeCharacter{03BB}{\ensuremath{\lambda}}
\DeclareUnicodeCharacter{03BC}{\ensuremath{\mu}}
\DeclareUnicodeCharacter{03BD}{\ensuremath{\nu}}
\DeclareUnicodeCharacter{03BE}{\ensuremath{\xi}}
\DeclareUnicodeCharacter{03C0}{\ensuremath{\pi}}
\DeclareUnicodeCharacter{03C1}{\ensuremath{\rho}}
\DeclareUnicodeCharacter{03C3}{\ensuremath{\sigma}}
\DeclareUnicodeCharacter{03C4}{\ensuremath{\tau}}
\DeclareUnicodeCharacter{03C5}{\ensuremath{\upsilon}}
\DeclareUnicodeCharacter{03C6}{\ensuremath{\varphi}}
\DeclareUnicodeCharacter{03C7}{\ensuremath{\chi}}
\DeclareUnicodeCharacter{03C8}{\ensuremath{\psi}}
\DeclareUnicodeCharacter{03C9}{\ensuremath{\omega}}
\DeclareUnicodeCharacter{0394}{\ensuremath{\Delta}}
\DeclareUnicodeCharacter{0398}{\ensuremath{\Theta}}
\DeclareUnicodeCharacter{039B}{\ensuremath{\Lambda}}
\DeclareUnicodeCharacter{03A0}{\ensuremath{\Pi}}
\DeclareUnicodeCharacter{03A3}{\ensuremath{\Sigma}}
\DeclareUnicodeCharacter{03A6}{\ensuremath{\Phi}}
\DeclareUnicodeCharacter{03A8}{\ensuremath{\Psi}}
\DeclareUnicodeCharacter{03A9}{\ensuremath{\Omega}}
\DeclareUnicodeCharacter{2200}{\ensuremath{\forall}}
\DeclareUnicodeCharacter{2203}{\ensuremath{\exists}}
\DeclareUnicodeCharacter{2207}{\ensuremath{\nabla}}
\DeclareUnicodeCharacter{2202}{\ensuremath{\partial}}
\DeclareUnicodeCharacter{2218}{\ensuremath{\circ}}
\DeclareUnicodeCharacter{2032}{\ensuremath{\prime}}
\DeclareUnicodeCharacter{2113}{\ensuremath{\ell}}
\DeclareUnicodeCharacter{210F}{\ensuremath{\hbar}}
\DeclareUnicodeCharacter{27E8}{\ensuremath{\langle}}
\DeclareUnicodeCharacter{27E9}{\ensuremath{\rangle}}
\DeclareUnicodeCharacter{2308}{\ensuremath{\lceil}}
\DeclareUnicodeCharacter{2309}{\ensuremath{\rceil}}
\DeclareUnicodeCharacter{230A}{\ensuremath{\lfloor}}
\DeclareUnicodeCharacter{230B}{\ensuremath{\rfloor}}
\DeclareUnicodeCharacter{222B}{\ensuremath{\int}}
\DeclareUnicodeCharacter{2211}{\ensuremath{\sum}}
\DeclareUnicodeCharacter{220F}{\ensuremath{\prod}}
\DeclareUnicodeCharacter{221E}{\ensuremath{\infty}}
\DeclareUnicodeCharacter{00B5}{\ensuremath{\mu}}
\DeclareUnicodeCharacter{2236}{\ensuremath{\colon}}
\DeclareUnicodeCharacter{2192}{\ensuremath{\to}}
\DeclareUnicodeCharacter{21A6}{\ensuremath{\mapsto}}
\DeclareUnicodeCharacter{21D2}{\ensuremath{\Rightarrow}}
\DeclareUnicodeCharacter{21D4}{\ensuremath{\Leftrightarrow}}
\DeclareUnicodeCharacter{2282}{\ensuremath{\subset}}
\DeclareUnicodeCharacter{2286}{\ensuremath{\subseteq}}
\DeclareUnicodeCharacter{2287}{\ensuremath{\supseteq}}
\DeclareUnicodeCharacter{222B}{\ensuremath{\int}}
\DeclareUnicodeCharacter{2205}{\ensuremath{\emptyset}}
\DeclareUnicodeCharacter{2299}{\ensuremath{\odot}}
\DeclareUnicodeCharacter{2220}{\ensuremath{\angle}}
\DeclareUnicodeCharacter{2135}{\ensuremath{\aleph}}
\DeclareUnicodeCharacter{207B}{\ensuremath{^{-}}}
\DeclareUnicodeCharacter{207A}{\ensuremath{^{+}}}
\DeclareUnicodeCharacter{220E}{\ensuremath{\blacksquare}}
\DeclareUnicodeCharacter{2061}{}
\DeclareUnicodeCharacter{2062}{}
\DeclareUnicodeCharacter{2063}{}
\DeclareUnicodeCharacter{2009}{\,}
\DeclareUnicodeCharacter{200A}{\,}
\DeclareUnicodeCharacter{200B}{}
\DeclareUnicodeCharacter{2212}{\ensuremath{-}}
\DeclareUnicodeCharacter{2194}{\ensuremath{\leftrightarrow}}
\DeclareUnicodeCharacter{2099}{\ensuremath{_{n}}}
\DeclareUnicodeCharacter{2098}{\ensuremath{_{m}}}
\DeclareUnicodeCharacter{2096}{\ensuremath{_{k}}}
\DeclareUnicodeCharacter{2095}{\ensuremath{_{h}}}
\DeclareUnicodeCharacter{2090}{\ensuremath{_{a}}}
\DeclareUnicodeCharacter{2091}{\ensuremath{_{e}}}
\DeclareUnicodeCharacter{2093}{\ensuremath{_{x}}}
\DeclareUnicodeCharacter{1D62}{\ensuremath{_{i}}}
\DeclareUnicodeCharacter{2C7C}{\ensuremath{_{j}}}
\DeclareUnicodeCharacter{2C7B}{\ensuremath{^{e}}}
\DeclareUnicodeCharacter{1D52}{\ensuremath{^{o}}}
\DeclareUnicodeCharacter{1D5}{\ensuremath{^{u}}}
\DeclareUnicodeCharacter{02E1}{\ensuremath{^{l}}}
\DeclareUnicodeCharacter{2295}{\ensuremath{\oplus}}
\DeclareUnicodeCharacter{27FA}{\ensuremath{\Longleftrightarrow}}
\DeclareUnicodeCharacter{27F9}{\ensuremath{\Longrightarrow}}
\DeclareUnicodeCharacter{2299}{\ensuremath{\odot}}
% Provide \tightlist used by pandoc-generated lists (no-op)
\providecommand{\tightlist}{\setlength{\itemsep}{0pt}\setlength{\parskip}{0pt}}
% pandoc longtable column specs: provide \real as passthrough and a
% `none` counter so that calc's \real{x} expansion path works.
\newcounter{none}
\providecommand{\real}[1]{#1}

% LISTINGS
\lstset{
    basicstyle=\small\ttfamily,
    breaklines=true, frame=single,
    keepspaces=true, showstringspaces=false,
    aboveskip=0.8em, belowskip=0.8em
}

% HEADER/FOOTER
\setlength{\headheight}{14pt}
\pagestyle{fancy}
\fancyhf{}
\fancyhead[R]{\thepage}
\renewcommand{\headrulewidth}{0.2pt}

% HYPERREF
\hypersetup{
    colorlinks=true,
    linkcolor=blue, citecolor=blue, urlcolor=blue,
    pdfauthor={Brieuc de La Fourniere}
}

% SPACING
\setstretch{1.2}
\setlength{\parskip}{0.4em}
\setlength{\parindent}{0pt}

% TITLE FORMATTING
\usepackage{titling}
\pretitle{\LARGE\bfseries}
\posttitle{\vspace{-0.4em}}
\preauthor{}
\postauthor{}
\predate{}
\postdate{}
\setlength{\droptitle}{-2.0em}

% CUSTOM COMMANDS — GIFT notation
\newcommand{\E}{\mathrm{E}}
\newcommand{\Gtwo}{\mathrm{G}_2}
\newcommand{\Kseven}{K_7}
\newcommand{\AdS}{\mathrm{AdS}}
\newcommand{\SU}{\mathrm{SU}}
\newcommand{\SO}{\mathrm{SO}}
\newcommand{\U}{\mathrm{U}}
\newcommand{\Sp}{\mathrm{Sp}}
\newcommand{\PSL}{\mathrm{PSL}}
\newcommand{\SM}{\mathrm{SM}}
\newcommand{\Tr}{\mathrm{Tr}}
\newcommand{\rank}{\mathrm{rank}}
\newcommand{\Vol}{\mathrm{Vol}}
\newcommand{\Hol}{\mathrm{Hol}}
\newcommand{\Ric}{\mathrm{Ric}}
\newcommand{\Rm}{\mathrm{Rm}}
\newcommand{\Scal}{\mathrm{Scal}}
\DeclareMathOperator{\diag}{diag}
\DeclareMathOperator{\cond}{cond}
\newcommand{\GIFT}{\textrm{GIFT}}
\newcommand{\CP}{\mathrm{CP}}
\newcommand{\CKM}{\mathrm{CKM}}
\newcommand{\PMNS}{\mathrm{PMNS}}
\newcommand{\EW}{\mathrm{EW}}
\newcommand{\Pl}{\mathrm{Pl}}
\newcommand{\DE}{\mathrm{DE}}
\newcommand{\DM}{\mathrm{DM}}
\newcommand{\KK}{\mathrm{KK}}
\newcommand{\NK}{\mathrm{NK}}
\newcommand{\TCS}{\mathrm{TCS}}
\newcommand{\PINN}{\mathrm{PINN}}
\newcommand{\btwo}{b_2}
\newcommand{\bthree}{b_3}
\newcommand{\typeI}{\textsc{I}}
\newcommand{\typeII}{\textsc{II}}
\newcommand{\typeIII}{\textsc{III}}
\newcommand{\typeIV}{\textsc{IV}}
\newcommand{\proven}{\textsc{Proven}}

% PDF bookmark safety
\pdfstringdefDisableCommands{%
  \def\textsubscript#1{#1}%
  \def\textsuperscript#1{#1}%
  \def\CP{CP}%
  \def\Gtwo{G2}%
  \def\Kseven{K7}%
  \def\GIFT{GIFT}%
  \def\E{E}%
  \def\SM{SM}%
}

% THEOREM ENVIRONMENTS
\newtheorem{theorem}{Theorem}[section]
\newtheorem{proposition}[theorem]{Proposition}
\newtheorem{lemma}[theorem]{Lemma}
\newtheorem{corollary}[theorem]{Corollary}
\theoremstyle{definition}
\newtheorem{definition}[theorem]{Definition}
\theoremstyle{remark}
\newtheorem{remark}[theorem]{Remark}

\fancyhead[L]{GIFT v3.4 --- Supplement S2}
\hypersetup{pdftitle={GIFT v3.4 Supplement S2: Complete Derivations}}

\title{%
\LARGE\textbf{Supplement S2: Complete Derivations}\\[0.3em]
\large Mathematical Proofs for All 33 Type I Predictions%
}
\author{Brieuc de La Fourni\`ere \\ \textit{Independent researcher, Beaune, France}}
\date{May 2026 \quad\textbar\quad v3.4}

\graphicspath{{../../../figures/}}

\begin{document}

\maketitle
\noindent\rule{\textwidth}{0.2pt}
\thispagestyle{fancy}

\tableofcontents
\newpage



\emph{This supplement provides mathematical derivations for the 33 Type I dimensionless predictions in the GIFT framework (Relations \#1--\#33, all Lean-certified). Each derivation proceeds from topological definitions to numerical predictions. Two additional BH remnant topological predictions (Relations \#34--\#35 in §24.3b) are included as an appendix; these are classified as Type IV structural diagnostics in the main 95-observable dataset.}

\textbf{Status}: All 33 Type I observables Lean-certified. 213 certificate conjuncts, 4 axioms, 0 sorry.

\textbf{Note on the 95-observable dataset}: The main paper presents 95 observables across 4 types (33 I + 19 II + 21 III + 22 IV). This supplement covers the 33 Type I predictions (direct algebraic from topology). Type II extensions (19 one-step physical extractions), Type III dynamical results (21 multi-step chains), and Type IV structural diagnostics (22 internal consistency checks, including BH remnant M\_res and N\_QNM from Pinčák et al.~2026 \cite{42}) are documented in Supplement S3 and the main text §4--§6.

The topological constants that determine these relations are described in S1.

\begin{center}\rule{0.5\linewidth}{0.5pt}\end{center}

\subsection{Table of Contents}\label{table-of-contents}

\begin{itemize}
\item
  \hyperref[part-i-foundations]{Part I: Foundations}
\item
  \hyperref[part-ii-foundational-theorems]{Part II: Foundational Theorems (4 relations)}
\item
  \hyperref[part-iii-gauge-sector]{Part III: Gauge Sector (2 relations)}
\item
  \hyperref[part-iv-lepton-sector]{Part IV: Lepton Sector (3 relations)}
\item
  \hyperref[part-v-quark-sector]{Part V: Quark Sector (1 relation)}
\item
  \hyperref[part-vi-neutrino-sector]{Part VI: Neutrino Sector (4 relations)}
\item
  \hyperref[part-vii-higgs-cosmology]{Part VII: Higgs \& Cosmology (3 relations)}
\item
  \hyperref[part-viii-summary-table]{Part VIII: Summary Table}
\item
  \hyperref[part-ix-observable-catalog]{Part IX: Observable Catalog}
\end{itemize}

\begin{center}\rule{0.5\linewidth}{0.5pt}\end{center}

\section{Part 0: Derivation Philosophy}\label{part-0-derivation-philosophy}

\subsection{0. What ``Derivation'' Means in GIFT}\label{what-derivation-means-in-gift}

Before presenting derivations, we clarify the logical structure:

\subsubsection{0.1 Inputs vs Outputs}\label{inputs-vs-outputs}

\textbf{Inputs} (taken as given):
- The octonion algebra $\mathbb{O}$ and its automorphism group G₂ = Aut($\mathbb{O}$)
- The E₈×E₈ gauge structure
- The K₇ manifold (compact G₂ 7-manifold with b₂ = 21, b₃ = 77; explicit metric certified in [A]; topological construction via Joyce-Karigiannis Z₂³ orbifold route (S1 §8.4); explicit closed-form neck ansatz with Wirtinger certificate (S1 §8.4, [D]))

\textbf{Outputs} (derived from inputs):
- The 33 Type I dimensionless predictions

\subsubsection{0.2 What We Do NOT Claim}\label{what-we-do-not-claim}

\begin{itemize}
\item
  That $\mathbb{O}$ → G₂ → K₇ is the unique geometry for physics
\item
  That the formulas are uniquely determined by geometric principles
\item
  That the selection rule for specific combinations (b₂/(b₃ + dim\_G₂) vs b₂/b₃) is understood
\end{itemize}

\subsubsection{0.3 What We Observe}\label{what-we-observe}

\begin{itemize}
\item
  Given the inputs, the outputs follow by algebra
\item
  The outputs match experiment to 0.73\% mean deviation (PDG 2024, NuFIT 6.0, 33 Type I core)
\item
  No continuous parameters are fitted
\end{itemize}

\subsubsection{0.4 Torsion Independence}\label{torsion-independence}

\textbf{Important}: All 33 Type I predictions use only topological invariants. The torsion T does not appear in any formula. Therefore:
- Predictions depend only on topology, not on the actual torsion value
- The value κ\_T = 1/61 is a topological bound, not a prediction ingredient

\subsubsection{0.5 Results from Analysis of the NK-Certified Metric}\label{results-from-analysis-of-the-nk-certified-metric}

Analysis of the NK-certified metric ([B]) computed all invariants directly from the explicit metric. Key results:

\begin{itemize}
\item
  \textbf{Topological realization}: The pair (b₂, b₃) = (21, 77) is realized topologically by the Joyce-Karigiannis Z₂³ orbifold route (S1 §8.4); an explicit closed-form neck ansatz with five-layer Wirtinger certificate is established in [D]. Orthogonal TCS is excluded by parity (b₂+b₃=98 even). Any Betti decomposition via Mayer-Vietoris is conditional on the building block identification, which is open.
\item
  \textbf{$\bar{\nu}$ = 0}: Lean-certified in TCSConstruction.lean. For rectangular TCS (k₊=k₋=1), $\bar{\nu}$=0 automatically (CGN invariant); this is conditional on TCS realization.
\item
  \textbf{New topological formula}: V\_min = √(Vol(K₇)/11), where 11 = b₃/n = 77/7. NK numerical: 219.90; formula: 221.24 (0.6\% ✓). Identity: V\_min² × b₃/n = Vol(K₇): a universal G₂ relation.
\item
  \textbf{U(1)² exact symmetry}: The T² directions θ and ψ are degenerate to 2×10⁻⁵ in the metric eigenvalues, propagating to S\_θ = S\_ψ = 6.1265 in all period integrals (exact to 2.6×10⁻⁸).
\item
  \textbf{SM gauge group origin}: NOT from ADE singularities (TCS is smooth). From g₂ ⊂ so(7) spectral structure and so(8) = g₂ ⊕ L ⊕ R triality → N\_gen = 3.
\end{itemize}

\begin{center}\rule{0.5\linewidth}{0.5pt}\end{center}

\section{Part I: Foundations}\label{part-i-foundations}

\subsection{1. Status Classification}\label{status-classification}

{\def\LTcaptype{none} % do not increment counter
\begin{longtable}[]{@{}
  >{\raggedright\arraybackslash}p{(\linewidth - 2\tabcolsep) * \real{0.4211}}
  >{\raggedright\arraybackslash}p{(\linewidth - 2\tabcolsep) * \real{0.5789}}@{}}
\toprule\noalign{}
\begin{minipage}[b]{\linewidth}\raggedright
Status
\end{minipage} & \begin{minipage}[b]{\linewidth}\raggedright
Criterion
\end{minipage} \\
\midrule\noalign{}
\endhead
\bottomrule\noalign{}
\endlastfoot
\textbf{VERIFIED} & Complete mathematical proof, exact result from topology \\
\textbf{VERIFIED (Lean 4)} & Verified by Lean 4 kernel with Mathlib (machine-checked) \\
\textbf{TOPOLOGICAL} & Direct consequence of manifold structure \\
\end{longtable}
}

\subsection{2. Notation}\label{notation}

{\def\LTcaptype{none} % do not increment counter
\begin{longtable}[]{@{}
  >{\raggedright\arraybackslash}p{(\linewidth - 4\tabcolsep) * \real{0.2963}}
  >{\raggedright\arraybackslash}p{(\linewidth - 4\tabcolsep) * \real{0.2593}}
  >{\raggedright\arraybackslash}p{(\linewidth - 4\tabcolsep) * \real{0.4444}}@{}}
\toprule\noalign{}
\begin{minipage}[b]{\linewidth}\raggedright
Symbol
\end{minipage} & \begin{minipage}[b]{\linewidth}\raggedright
Value
\end{minipage} & \begin{minipage}[b]{\linewidth}\raggedright
Definition
\end{minipage} \\
\midrule\noalign{}
\endhead
\bottomrule\noalign{}
\endlastfoot
dim(E₈) & 248 & E₈ Lie algebra dimension \\
rank(E₈) & 8 & E₈ Cartan subalgebra dimension \\
dim(G₂) & 14 & G₂ holonomy group dimension \\
dim(K₇) & 7 & Internal manifold dimension \\
b₂(K₇) & 21 & Second Betti number \\
b₃(K₇) & 77 & Third Betti number \\
H* & 99 & Effective cohomology = b₂ + b₃ + 1 \\
dim(J₃(O)) & 27 & Exceptional Jordan algebra dimension \\
N\_gen & 3 & Number of fermion generations \\
p₂ & 2 & Binary duality parameter \\
Weyl & 5 & Weyl factor: (dim(G₂)+1)/N\_gen = b₂/N\_gen - p₂ = dim(G₂) - rank(E₈) - 1 \\
\end{longtable}
}

\begin{center}\rule{0.5\linewidth}{0.5pt}\end{center}

\section{Part II: Foundational Theorems}\label{part-ii-foundational-theorems}

\subsection{3. Relation \#1: Generation Number N\_gen = 3}\label{relation-1-generation-number-n_gen-3}

\textbf{Statement}: The number of fermion generations is exactly 3.

\textbf{Classification}: VERIFIED (three independent derivations)

\subsubsection{Proof Method 1: Fundamental Topological Constraint}\label{proof-method-1-fundamental-topological-constraint}

\emph{Theorem}: For G₂ holonomy manifold K₇ with E₈ gauge structure:

\[(\text{rank}(E_8) + N_{\text{gen}}) \cdot b_2(K_7) = N_{\text{gen}} \cdot b_3(K_7)\]

\emph{Derivation}:
\[(8 + N_{\text{gen}}) \times 21 = N_{\text{gen}} \times 77\]
\[168 + 21 \cdot N_{\text{gen}} = 77 \cdot N_{\text{gen}}\]
\[168 = 56 \cdot N_{\text{gen}}\]
\[N_{\text{gen}} = \frac{168}{56} = 3\]

\emph{Verification}:
- LHS: (8 + 3) × 21 = 231
- RHS: 3 × 77 = 231 ✓

\subsubsection{Proof Method 2: Atiyah-Singer Index Theorem}\label{proof-method-2-atiyah-singer-index-theorem}

\[\text{Index}(D_A) = \left( 77 - \frac{8}{3} \times 21 \right) \times \frac{1}{7} = 3\]

\textbf{Status}: VERIFIED ∎

\begin{center}\rule{0.5\linewidth}{0.5pt}\end{center}

\subsection{4. Relation \#2: Hierarchy Parameter τ = 3472/891}\label{relation-2-hierarchy-parameter-ux3c4-3472891}

\textbf{Statement}: The hierarchy parameter is exactly rational.

\textbf{Classification}: VERIFIED

\subsubsection{Proof}\label{proof}

\emph{Step 1: Definition from topological integers}
\[\tau := \frac{\dim(E_8 \times E_8) \cdot b_2(K_7)}{\dim(J_3(\mathbb{O})) \cdot H^*}\]

\emph{Step 2: Substitute values}
\[\tau = \frac{496 \times 21}{27 \times 99} = \frac{10416}{2673}\]

\emph{Step 3: Reduce}
\[\gcd(10416, 2673) = 3\]
\[\tau = \frac{3472}{891}\]

\emph{Step 4: Prime factorization}
\[\tau = \frac{2^4 \times 7 \times 31}{3^4 \times 11}\]

\emph{Step 5: Numerical value}
\[\tau = 3.8967452300785634...\]

\textbf{Status}: VERIFIED ∎

\begin{center}\rule{0.5\linewidth}{0.5pt}\end{center}

\subsection{5. Relation \#3: Torsion Capacity κ\_T = 1/61}\label{relation-3-torsion-capacity-ux3ba_t-161}

\textbf{Statement}: The topological torsion capacity equals exactly 1/61.

\textbf{Classification}: TOPOLOGICAL (structural parameter, not physical prediction)

\subsubsection{Proof}\label{proof-1}

\emph{Step 1: Define from cohomology}
\[61 = b_3(K_7) - \dim(G_2) - p_2 = 77 - 14 - 2 = 61\]

\emph{Step 2: Formula}
\[\kappa_T = \frac{1}{b_3 - \dim(G_2) - p_2} = \frac{1}{61}\]

\emph{Step 3: Geometric interpretation}
- 61 = effective degrees of freedom available for torsional deformation
- 61 = dim(F₄) + N\_gen² = 52 + 9

\subsubsection{Clarification}\label{clarification}

{\def\LTcaptype{none} % do not increment counter
\begin{longtable}[]{@{}
  >{\raggedright\arraybackslash}p{(\linewidth - 4\tabcolsep) * \real{0.3448}}
  >{\raggedright\arraybackslash}p{(\linewidth - 4\tabcolsep) * \real{0.4138}}
  >{\raggedright\arraybackslash}p{(\linewidth - 4\tabcolsep) * \real{0.2414}}@{}}
\toprule\noalign{}
\begin{minipage}[b]{\linewidth}\raggedright
Quantity
\end{minipage} & \begin{minipage}[b]{\linewidth}\raggedright
Definition
\end{minipage} & \begin{minipage}[b]{\linewidth}\raggedright
Value
\end{minipage} \\
\midrule\noalign{}
\endhead
\bottomrule\noalign{}
\endlastfoot
κ\_T & Topological capacity & 1/61 (fixed) \\
T\_base & Torsion for torsion-free metric (Joyce) & \textbf{0} (by theorem) \\
T\_physical & Effective torsion for interactions & \textbf{Open question} \\
\end{longtable}
}

\textbf{Role in predictions}: κ\_T appears in only one formula (α⁻¹, as a small correction term det(g)×κ\_T ≈ 0.033). The other 32 Type I predictions are independent of torsion capacity. It is primarily a structural parameter characterizing K₇, not a directly measured observable.

\textbf{Joyce's theorem}: For compact G₂ manifolds satisfying perturbation bounds, guarantees existence of a torsion-free metric. The Newton-Kantorovich certificate (Paper I) establishes these bounds for K₇.

\textbf{Status}: TOPOLOGICAL (structural, not predictive) ∎

\begin{center}\rule{0.5\linewidth}{0.5pt}\end{center}

\subsection{6. Relation \#4: Metric Determinant det(g) = 65/32}\label{relation-4-metric-determinant-detg-6532}

\textbf{Statement}: The K₇ metric determinant is exactly 65/32.

\textbf{Classification}: METRIC NORMALIZATION (not TOPOLOGICAL, see S1 §10.3)

\subsubsection{Proof}\label{proof-2}

\emph{Step 1: Define from structural normalization integers}
\[\det(g) = p_2 + \frac{1}{b_2 + \dim(G_2) - N_{gen}}\]

\emph{Step 2: Compute denominator}
\[b_2 + \dim(G_2) - N_{gen} = 21 + 14 - 3 = 32\]

\emph{Step 3: Compute determinant}
\[\det(g) = 2 + \frac{1}{32} = \frac{65}{32}\]

\emph{Step 4: Alternative derivation}
\[\det(g) = \frac{\text{Weyl} \times (\text{rank}(E_8) + \text{Weyl})}{2^5} = \frac{5 \times 13}{32} = \frac{65}{32}\]

\textbf{Verification}: The analytical metric g = (65/32)\^{}\{1/7\} × I₇ has det(g) = {[}(65/32)\^{}\{1/7\}{]}\^{}7 = 65/32 exactly, consistent with the topological formula.

\textbf{Epistemic note}: Three independent algebraic paths converge to 65/32, which is suggestive but does not constitute a derivation from topology: the metric optimization was constrained to this normalization target, and the formulas were identified post-hoc. See S1 §10.3 for full discussion. Six observables depending on det(g)\_num or det(g)\_den carry this normalization dependence.

\textbf{Status}: STRUCTURAL (metric normalization, algebraically exact; see S1 §10.3) ∎

\begin{center}\rule{0.5\linewidth}{0.5pt}\end{center}

\section{Part III: Gauge Sector}\label{part-iii-gauge-sector}

\subsection{7. Relation \#5: Weinberg Angle sin²θ\_W = 3/13}\label{relation-5-weinberg-angle-sinuxb2ux3b8_w-313}

\textbf{Statement}: The weak mixing angle has exact rational form 3/13.

\textbf{Classification}: VERIFIED

\subsubsection{Proof}\label{proof-3}

\emph{Step 1: Define ratio from Betti numbers}
\[\sin^2\theta_W = \frac{b_2(K_7)}{b_3(K_7) + \dim(G_2)} = \frac{21}{77 + 14} = \frac{21}{91}\]

\emph{Step 2: Simplify}
\[\gcd(21, 91) = 7\]
\[\sin^2\theta_W = \frac{3}{13} = 0.230769...\]

\emph{Step 3: Experimental comparison}

{\def\LTcaptype{none} % do not increment counter
\begin{longtable}[]{@{}ll@{}}
\toprule\noalign{}
Quantity & Value \\
\midrule\noalign{}
\endhead
\bottomrule\noalign{}
\endlastfoot
Experimental (PDG 2024) & 0.23122 ± 0.00004 \\
GIFT prediction & 0.230769 \\
Deviation & 0.195\% \\
\end{longtable}
}

\textbf{Status}: VERIFIED ∎

\begin{center}\rule{0.5\linewidth}{0.5pt}\end{center}

\subsection{8. Relation \#6: Strong Coupling α\_s = √2/12}\label{relation-6-strong-coupling-ux3b1_s-212}

\textbf{Statement}: The strong coupling at M\_Z scale.

\textbf{Classification}: TOPOLOGICAL

\subsubsection{Proof}\label{proof-4}

\emph{Formula}:
\[\alpha_s(M_Z) = \frac{\sqrt{2}}{\dim(G_2) - p_2} = \frac{\sqrt{2}}{14 - 2} = \frac{\sqrt{2}}{12}\]

\emph{Components}:
- √2: E₈ root length
- 12 = dim(G₂) - p₂: Effective gauge degrees of freedom

\emph{Numerical value}: α\_s = 0.117851

\emph{Experimental comparison}:

{\def\LTcaptype{none} % do not increment counter
\begin{longtable}[]{@{}ll@{}}
\toprule\noalign{}
Quantity & Value \\
\midrule\noalign{}
\endhead
\bottomrule\noalign{}
\endlastfoot
Experimental & 0.1180 ± 0.0009 \\
GIFT prediction & 0.11785 \\
Deviation & 0.127\% \\
\end{longtable}
}

\textbf{Status}: TOPOLOGICAL ∎

\begin{center}\rule{0.5\linewidth}{0.5pt}\end{center}

\section{Part IV: Lepton Sector}\label{part-iv-lepton-sector}

\subsection{9. Relation \#7: Koide Parameter Q = 2/3}\label{relation-7-koide-parameter-q-23}

\textbf{Statement}: The Koide parameter equals exactly 2/3.

\textbf{Classification}: VERIFIED

\subsubsection{Proof}\label{proof-5}

\emph{Formula}:
\[Q_{\text{Koide}} = \frac{\dim(G_2)}{b_2(K_7)} = \frac{14}{21} = \frac{2}{3}\]

\emph{Physical definition}:
\[Q = \frac{m_e + m_\mu + m_\tau}{(\sqrt{m_e} + \sqrt{m_\mu} + \sqrt{m_\tau})^2}\]

\emph{Experimental comparison}:

{\def\LTcaptype{none} % do not increment counter
\begin{longtable}[]{@{}ll@{}}
\toprule\noalign{}
Quantity & Value \\
\midrule\noalign{}
\endhead
\bottomrule\noalign{}
\endlastfoot
Experimental & 0.666661 ± 0.000007 \\
GIFT prediction & 0.666667 \\
Deviation & 0.0009\% \\
\end{longtable}
}

\textbf{Status}: VERIFIED ∎

\begin{center}\rule{0.5\linewidth}{0.5pt}\end{center}

\subsection{10. Relation \#8: Tau-Electron Mass Ratio m\_τ/m\_e = 3477}\label{relation-8-tau-electron-mass-ratio-m_ux3c4m_e-3477}

\textbf{Statement}: The tau-electron mass ratio is exactly 3477.

\textbf{Classification}: VERIFIED

\subsubsection{Proof}\label{proof-6}

\emph{Formula}:
\[\frac{m_\tau}{m_e} = \dim(K_7) + 10 \cdot \dim(E_8) + 10 \cdot H^*\]
\[= 7 + 10 \times 248 + 10 \times 99 = 7 + 2480 + 990 = 3477\]

\emph{Prime factorization}:
\[3477 = 3 \times 19 \times 61 = N_{gen} \times prime(8) \times \kappa_T^{-1}\]

\emph{Period integral cross-confirmation}: Analysis of the NK-certified metric ([B], §5) extracts associative 3-cycle volumes directly. The SD associative volumes yield ln(m\_τ/m\_e) = 8.154 from the Fano-plane volume hierarchy, giving e\^{}8.154 = 3477 ✓. This provides an independent geometric confirmation of the algebraic prediction.

\emph{Experimental comparison}:

{\def\LTcaptype{none} % do not increment counter
\begin{longtable}[]{@{}ll@{}}
\toprule\noalign{}
Quantity & Value \\
\midrule\noalign{}
\endhead
\bottomrule\noalign{}
\endlastfoot
Experimental & 3477.15 ± 0.05 \\
GIFT prediction & 3477 (exact) \\
Deviation & 0.0043\% \\
\end{longtable}
}

\textbf{Status}: VERIFIED ∎

\begin{center}\rule{0.5\linewidth}{0.5pt}\end{center}

\subsection{11. Relation \#9: Muon-Electron Mass Ratio}\label{relation-9-muon-electron-mass-ratio}

\textbf{Statement}: m\_μ/m\_e = 27\^{}φ

\textbf{Classification}: TOPOLOGICAL

\subsubsection{Proof}\label{proof-7}

\emph{Formula}:
\[\frac{m_\mu}{m_e} = [\dim(J_3(\mathbb{O}))]^\phi = 27^\phi = 207.012\]

\emph{Components}:
- 27 = dim(J₃(O)): Exceptional Jordan algebra
- φ = (1+√5)/2: Golden ratio from McKay correspondence E₈ ↔ 2I (binary icosahedral group)

\emph{Experimental comparison}:

{\def\LTcaptype{none} % do not increment counter
\begin{longtable}[]{@{}ll@{}}
\toprule\noalign{}
Quantity & Value \\
\midrule\noalign{}
\endhead
\bottomrule\noalign{}
\endlastfoot
Experimental & 206.768 \\
GIFT prediction & 207.01 \\
Deviation & 0.1179\% \\
\end{longtable}
}

\textbf{Note on φ}: The golden ratio is the only non-integer constant among the 33 Type I predictions and does not appear in the 20 structural constants (S3.3). Its derivation from E₈ via the McKay correspondence is well-established mathematically but constitutes an additional identification step beyond the direct topological reasoning used for the other 32 predictions.

\textbf{Status}: TOPOLOGICAL (with caveat: φ external to the 20 structural constants) ∎

\begin{center}\rule{0.5\linewidth}{0.5pt}\end{center}

\section{Part V: Quark Sector}\label{part-v-quark-sector}

\subsection{12. Relation \#10: Strange-Down Ratio m\_s/m\_d = 20}\label{relation-10-strange-down-ratio-m_sm_d-20}

\textbf{Statement}: The strange-down quark mass ratio is exactly 20.

\textbf{Classification}: VERIFIED

\subsubsection{Proof}\label{proof-8}

\emph{Formula}:
\[\frac{m_s}{m_d} = p_2^2 \times \text{Weyl} = 4 \times 5 = 20\]

\emph{Geometric interpretation}:
- p₂² = 4: Binary structure squared
- Weyl = 5: Pentagonal symmetry

\emph{Experimental comparison}:

{\def\LTcaptype{none} % do not increment counter
\begin{longtable}[]{@{}ll@{}}
\toprule\noalign{}
Quantity & Value \\
\midrule\noalign{}
\endhead
\bottomrule\noalign{}
\endlastfoot
Experimental & 20.0 ± 1.0 \\
GIFT prediction & 20 (exact) \\
Deviation & 0.00\% \\
\end{longtable}
}

\textbf{Status}: VERIFIED ∎

\begin{center}\rule{0.5\linewidth}{0.5pt}\end{center}

\subsection{12b. Relation \#10b: Charm-Strange Ratio m\_c/m\_s = 246/21}\label{b.-relation-10b-charm-strange-ratio-m_cm_s-24621}

\textbf{Statement}: The charm-strange quark mass ratio.

\textbf{Classification}: TOPOLOGICAL

\subsubsection{Proof}\label{proof-9}

\emph{Formula}:
\[\frac{m_c}{m_s} = \frac{\dim(E_8) - p_2}{b_2(K_7)} = \frac{248 - 2}{21} = \frac{246}{21} = 11.714...\]

\emph{Components}:
- 246 = dim(E₈) - p₂: Effective E₈ dimension
- 21 = b₂(K₇): Second Betti number

\emph{Experimental comparison}:

{\def\LTcaptype{none} % do not increment counter
\begin{longtable}[]{@{}ll@{}}
\toprule\noalign{}
Quantity & Value \\
\midrule\noalign{}
\endhead
\bottomrule\noalign{}
\endlastfoot
Experimental & 11.7 ± 0.3 \\
GIFT prediction & 11.714 \\
Deviation & 0.12\% \\
\end{longtable}
}

\emph{TCS-level alternative}: The charm-strange ratio also admits a decomposition at the TCS building-block level:

\[\frac{m_c}{m_s} = b_2(M_1) + \frac{\dim(K_7)}{b_2(M_2)} = 11 + \frac{7}{10} = \frac{117}{10} = 11.700\]

This additive form reflects the adiabatic decomposition of the KK mass operator across the two TCS halves. The M₁ contribution (b₂(M₁) = 11) dominates, with a geometric correction from M₂ (dim(K₇)/b₂(M₂) = 7/10). Both formulas are within 0.05σ of the experimental central value.

\textbf{Status}: TOPOLOGICAL ∎

\begin{center}\rule{0.5\linewidth}{0.5pt}\end{center}

\subsection{12c. Relation \#10c: Bottom-Top Ratio m\_b/m\_t = 1/42}\label{c.-relation-10c-bottom-top-ratio-m_bm_t-142}

\textbf{Statement}: The bottom-top quark mass ratio involves the constant 42 = p₂ × N\_gen × dim(K₇).

\textbf{Classification}: TOPOLOGICAL

\subsubsection{Proof}\label{proof-10}

\emph{Step 1: Define the structural constant}
\[42 = p_2 \times N_{gen} \times \dim(K_7) = 2 \times 3 \times 7\]

This constant 42 also equals 2 × b₂ = 2 × 21.

\emph{Step 2: Formula}
\[\frac{m_b}{m_t} = \frac{b_0}{42} = \frac{1}{42} = 0.02381...\]

\emph{Components}:
- b₀ = 1: Zeroth Betti number
- 42: Structural constant from K₇ geometry

\emph{Experimental comparison}:

{\def\LTcaptype{none} % do not increment counter
\begin{longtable}[]{@{}ll@{}}
\toprule\noalign{}
Quantity & Value \\
\midrule\noalign{}
\endhead
\bottomrule\noalign{}
\endlastfoot
Experimental & 0.024 ± 0.001 \\
GIFT prediction & 0.02381 \\
Deviation & 0.79\% \\
\end{longtable}
}

\emph{Geometric interpretation}: The same constant 42 appears in the cosmological ratio Ω\_DM/Ω\_b = (1 + 42)/8 = 43/8 (Section 18b), connecting quark physics to cosmological structure through the K₇ geometry.

\textbf{Status}: TOPOLOGICAL ∎

\begin{center}\rule{0.5\linewidth}{0.5pt}\end{center}

\subsection{12d. Relation \#10d: Up-Down Ratio m\_u/m\_d = 79/168}\label{d.-relation-10d-up-down-ratio-m_um_d-79168}

\textbf{Statement}: The up-down quark mass ratio.

\textbf{Classification}: TOPOLOGICAL

\subsubsection{Proof}\label{proof-11}

\emph{Formula}:
\[\frac{m_u}{m_d} = \frac{b_0 + \dim(E_6)}{|PSL_2(7)|} = \frac{1 + 78}{168} = \frac{79}{168} = 0.4702...\]

\emph{Components}:
- dim(E₆) = 78: Exceptional Lie algebra dimension
- \textbar PSL₂(7)\textbar{} = 168: Order of the simple group PSL₂(7) = rank(E₈) × b₂

\emph{Experimental comparison}:

{\def\LTcaptype{none} % do not increment counter
\begin{longtable}[]{@{}ll@{}}
\toprule\noalign{}
Quantity & Value \\
\midrule\noalign{}
\endhead
\bottomrule\noalign{}
\endlastfoot
Experimental & 0.47 ± 0.03 \\
GIFT prediction & 0.4702 \\
Deviation & 0.05\% \\
\end{longtable}
}

\textbf{Status}: TOPOLOGICAL ∎

\begin{center}\rule{0.5\linewidth}{0.5pt}\end{center}

\section{Part V-B: CKM Matrix}\label{part-v-b-ckm-matrix}

\subsection{12e. Relation \#10e: Cabibbo Angle sin²θ₁₂(CKM) = 7/31}\label{e.-relation-10e-cabibbo-angle-sinuxb2ux3b8ux2081ux2082ckm-731}

\textbf{Statement}: The CKM Cabibbo mixing angle.

\textbf{Classification}: TOPOLOGICAL

\subsubsection{Proof}\label{proof-12}

\emph{Formula}:
\[\sin^2\theta_{12}^{CKM} = \frac{\dim(\text{fund}_{E_7})}{\dim(E_8)} = \frac{56}{248} = \frac{7}{31} = 0.2258...\]

\emph{Alternative expressions}:
- (b₃ - b₂)/dim(E₈) = (77 - 21)/248 = 56/248
- (2b₂ + dim(G₂))/dim(E₈) = (42 + 14)/248 = 56/248

\emph{Experimental comparison}:

{\def\LTcaptype{none} % do not increment counter
\begin{longtable}[]{@{}ll@{}}
\toprule\noalign{}
Quantity & Value \\
\midrule\noalign{}
\endhead
\bottomrule\noalign{}
\endlastfoot
Experimental & 0.2250 ± 0.0006 \\
GIFT prediction & 0.2258 \\
Deviation & 0.36\% \\
\end{longtable}
}

\textbf{Status}: TOPOLOGICAL ∎

\begin{center}\rule{0.5\linewidth}{0.5pt}\end{center}

\subsection{12f. Relation \#10f: Wolfenstein A Parameter = 83/99}\label{f.-relation-10f-wolfenstein-a-parameter-8399}

\textbf{Statement}: The Wolfenstein A parameter of the CKM matrix.

\textbf{Classification}: TOPOLOGICAL

\subsubsection{Proof}\label{proof-13}

\emph{Formula}:
\[A_{\text{Wolf}} = \frac{\text{Weyl} + \dim(E_6)}{H^*} = \frac{5 + 78}{99} = \frac{83}{99} = 0.8384...\]

\emph{Alternative expression}:
- (b₃ + p₂ × N\_gen)/H* = (77 + 6)/99 = 83/99

\emph{Experimental comparison}:

{\def\LTcaptype{none} % do not increment counter
\begin{longtable}[]{@{}ll@{}}
\toprule\noalign{}
Quantity & Value \\
\midrule\noalign{}
\endhead
\bottomrule\noalign{}
\endlastfoot
Experimental & 0.836 ± 0.015 \\
GIFT prediction & 0.8384 \\
Deviation & 0.29\% \\
\end{longtable}
}

\textbf{Status}: TOPOLOGICAL ∎

\begin{center}\rule{0.5\linewidth}{0.5pt}\end{center}

\subsection{12g. Relation \#10g: CKM θ₂₃ Mixing sin²θ₂₃(CKM) = 1/24}\label{g.-relation-10g-ckm-ux3b8ux2082ux2083-mixing-sinuxb2ux3b8ux2082ux2083ckm-124}

\textbf{Statement}: The CKM 23-mixing angle.

\textbf{Classification}: TOPOLOGICAL

\subsubsection{Proof}\label{proof-14}

\emph{Formula}:
\[\sin^2\theta_{23}^{CKM} = \frac{\dim(K_7)}{|PSL_2(7)|} = \frac{7}{168} = \frac{1}{24} = 0.04167...\]

\emph{Experimental comparison}:

{\def\LTcaptype{none} % do not increment counter
\begin{longtable}[]{@{}ll@{}}
\toprule\noalign{}
Quantity & Value \\
\midrule\noalign{}
\endhead
\bottomrule\noalign{}
\endlastfoot
Experimental & 0.0412 ± 0.0008 \\
GIFT prediction & 0.04167 \\
Deviation & 1.13\% \\
\end{longtable}
}

\textbf{Status}: TOPOLOGICAL ∎

\begin{center}\rule{0.5\linewidth}{0.5pt}\end{center}

\section{Part VI: Neutrino Sector}\label{part-vi-neutrino-sector}

\subsection{13. Relation \#11: CP Violation Phase δ\_CP = 197°}\label{relation-11-cp-violation-phase-ux3b4_cp-197}

\textbf{Statement}: The CP violation phase is exactly 197°.

\textbf{Classification}: VERIFIED (with NuFIT 6.0 tension discussed in Appendix F)

\subsubsection{Proof}\label{proof-15}

\emph{Formula}:
\[\delta_{CP} = \dim(K_7) \cdot \dim(G_2) + H^* = 7 \times 14 + 99 = 98 + 99 = 197°\]

\emph{Experimental comparison}:

{\def\LTcaptype{none} % do not increment counter
\begin{longtable}[]{@{}ll@{}}
\toprule\noalign{}
Quantity & Value \\
\midrule\noalign{}
\endhead
\bottomrule\noalign{}
\endlastfoot
Experimental (NuFIT 5.2, 2022) & 197° ± 24° \\
Experimental (NuFIT 6.0, Oct 2024) & 177° ± 20° \\
GIFT prediction & 197° (exact) \\
Deviation (vs NuFIT 5.2) & 0.00\% \\
Deviation (vs NuFIT 6.0) & 11.3\% (within 1σ) \\
\end{longtable}
}

\textbf{NuFIT 6.0 tension}: The central value shifted from 197° to 177° between NuFIT 5.2
and 6.0, but the NuFIT collaboration flags δ\_CP as one of the least constrained
parameters, with substantial dependence on reactor flux assumptions. The value 197°
lies at the 1σ boundary (177+20=197). DUNE (2028--2040) will resolve this definitively.

\textbf{The canonical GIFT prediction remains δ\_CP = 197°.} A potential compactification
correction (62/69 factor, yielding 177.01°) is documented in Appendix F as contingent
on experimental confirmation. See also §4.4.1 of the main paper.

\textbf{Status}: VERIFIED, original prediction, pending DUNE ∎

\begin{center}\rule{0.5\linewidth}{0.5pt}\end{center}

\subsection{14. Relation \#12: Reactor Mixing Angle θ₁₃ = π/21}\label{relation-12-reactor-mixing-angle-ux3b8ux2081ux2083-ux3c021}

\textbf{Statement}: The reactor neutrino mixing angle.

\textbf{Classification}: TOPOLOGICAL

\subsubsection{Proof}\label{proof-16}

\emph{Formula}:
\[\theta_{13} = \frac{\pi}{b_2(K_7)} = \frac{\pi}{21} = 8.571°\]

\emph{Experimental comparison}:

{\def\LTcaptype{none} % do not increment counter
\begin{longtable}[]{@{}ll@{}}
\toprule\noalign{}
Quantity & Value \\
\midrule\noalign{}
\endhead
\bottomrule\noalign{}
\endlastfoot
Experimental (NuFIT 6.0) & 8.54° ± 0.12° \\
GIFT prediction & 8.571° \\
Deviation & 0.37\% \\
\end{longtable}
}

\emph{TCS-level refinement}: At the building-block level, the reactor angle admits a purely rational formula:

\[\theta_{13} = \frac{\dim(E_8) - \dim(G_2) \cdot b_2(M_1)}{b_2(M_1)} = \frac{248 - 14 \times 11}{11} = \frac{94}{11} = 8.545°\]

The numerator 94 = dim(E₈) − dim(G₂)×b₂(M₁) is the residual E₈ gauge content after subtracting the G₂ holonomy contribution to M₁. This formula is 5.8× more precise (0.064\% vs 0.37\%) and uses only algebraic invariants (no π), but works in degrees rather than radians.

\textbf{Status}: TOPOLOGICAL ∎

\begin{center}\rule{0.5\linewidth}{0.5pt}\end{center}

\subsection{15. Relation \#13: Atmospheric Mixing Angle θ₂₃}\label{relation-13-atmospheric-mixing-angle-ux3b8ux2082ux2083}

\textbf{Statement}: The atmospheric neutrino mixing angle.

\textbf{Classification}: TOPOLOGICAL

\subsubsection{Proof}\label{proof-17}

\emph{Formula}:
\[\theta_{23} = \arcsin\left(\frac{b_3 - p_2}{H^*}\right) = \arcsin\left(\frac{75}{99}\right) = \arcsin\left(\frac{25}{33}\right) = 49.251°\]

\emph{Components}:
- b₃ = 77: Third Betti number (3-cycles of K₇)
- p₂ = 2: Pontryagin class contribution (spin structure correction)
- H* = 99: Effective cohomology (b₂ + b₃ + 1)

\emph{Physical interpretation}:
The atmospheric mixing angle θ₂₃ governs τ-μ flavor mixing. The formula (b₃ - p₂)/H* represents the relative weight of spin-corrected 3-cycles in the total cohomology. This captures how the τ-μ sector couples through the 3-cycle topology of K₇, with the Pontryagin correction accounting for the spin structure that distinguishes fermionic generations.

\emph{Experimental comparison}:

{\def\LTcaptype{none} % do not increment counter
\begin{longtable}[]{@{}ll@{}}
\toprule\noalign{}
Quantity & Value \\
\midrule\noalign{}
\endhead
\bottomrule\noalign{}
\endlastfoot
Experimental (NuFIT 5.3) & 49.3° ± 1.0° \\
GIFT prediction & 49.251° \\
Deviation & 0.10\% \\
\end{longtable}
}

\textbf{Status}: TOPOLOGICAL ∎

\begin{center}\rule{0.5\linewidth}{0.5pt}\end{center}

\subsection{16. Relation \#14: Solar Mixing Angle θ₁₂}\label{relation-14-solar-mixing-angle-ux3b8ux2081ux2082}

\textbf{Statement}: The solar neutrino mixing angle.

\textbf{Classification}: TOPOLOGICAL

\subsubsection{Proof}\label{proof-18}

\emph{Formula}:
\[\theta_{12} = \arctan\left(\sqrt{\frac{\delta}{\gamma_{\text{GIFT}}}}\right) = 33.419°\]

\emph{Components}:
- δ = 2π/Weyl² = 2π/25
- γ\_GIFT = 511/884

\emph{Derivation of γ\_GIFT}:
\[\gamma_{\text{GIFT}} = \frac{2 \cdot \text{rank}(E_8) + 5 \cdot H^*}{10 \cdot \dim(G_2) + 3 \cdot \dim(E_8)} = \frac{511}{884}\]

\emph{Experimental comparison}:

{\def\LTcaptype{none} % do not increment counter
\begin{longtable}[]{@{}ll@{}}
\toprule\noalign{}
Quantity & Value \\
\midrule\noalign{}
\endhead
\bottomrule\noalign{}
\endlastfoot
Experimental (NuFIT 5.3) & 33.41° ± 0.75° \\
GIFT prediction & 33.40° \\
Deviation & 0.030\% \\
\end{longtable}
}

\textbf{Status}: TOPOLOGICAL ∎

\begin{center}\rule{0.5\linewidth}{0.5pt}\end{center}

\subsection{16b. PMNS Matrix: sin² Form}\label{b.-pmns-matrix-sinuxb2-form}

The PMNS mixing angles can also be expressed directly as sin² values, providing alternative topological formulas.

\subsubsection{Relation \#14b: sin²θ₁₂(PMNS) = 4/13}\label{relation-14b-sinuxb2ux3b8ux2081ux2082pmns-413}

\textbf{Formula}:
\[\sin^2\theta_{12}^{PMNS} = \frac{b_0 + N_{gen}}{\alpha_{sum}} = \frac{1 + 3}{13} = \frac{4}{13} = 0.3077...\]

\emph{Components}:
- α\_sum = 13: Anomaly coefficient sum
- b₀ + N\_gen = 4: Cohomological + generation count

{\def\LTcaptype{none} % do not increment counter
\begin{longtable}[]{@{}ll@{}}
\toprule\noalign{}
Quantity & Value \\
\midrule\noalign{}
\endhead
\bottomrule\noalign{}
\endlastfoot
Experimental & 0.307 ± 0.013 \\
GIFT prediction & 0.3077 \\
Deviation & 0.23\% \\
\end{longtable}
}

\subsubsection{Relation \#14c: sin²θ₂₃(PMNS) = 6/11}\label{relation-14c-sinuxb2ux3b8ux2082ux2083pmns-611}

\textbf{Formula}:
\[\sin^2\theta_{23}^{PMNS} = \frac{D_{bulk} - \text{Weyl}}{D_{bulk}} = \frac{11 - 5}{11} = \frac{6}{11} = 0.5455...\]

\emph{Alternative expression}:
- 42/b₃ = 42/77 = 6/11 (after reduction)

{\def\LTcaptype{none} % do not increment counter
\begin{longtable}[]{@{}ll@{}}
\toprule\noalign{}
Quantity & Value \\
\midrule\noalign{}
\endhead
\bottomrule\noalign{}
\endlastfoot
Experimental & 0.546 ± 0.021 \\
GIFT prediction & 0.5455 \\
Deviation & 0.10\% \\
\end{longtable}
}

\subsubsection{Relation \#14d: sin²θ₁₃(PMNS) = 11/496}\label{relation-14d-sinuxb2ux3b8ux2081ux2083pmns-11496}

\textbf{Formula}:
\[\sin^2\theta_{13}^{PMNS} = \frac{D_{bulk}}{\dim(E_8 \times E_8)} = \frac{11}{496} = 0.02218...\]

{\def\LTcaptype{none} % do not increment counter
\begin{longtable}[]{@{}ll@{}}
\toprule\noalign{}
Quantity & Value \\
\midrule\noalign{}
\endhead
\bottomrule\noalign{}
\endlastfoot
Experimental & 0.0220 ± 0.0007 \\
GIFT prediction & 0.02218 \\
Deviation & 0.81\% \\
\end{longtable}
}

\textbf{Status}: TOPOLOGICAL ∎

\begin{center}\rule{0.5\linewidth}{0.5pt}\end{center}

\section{Part VII: Higgs \& Cosmology}\label{part-vii-higgs-cosmology}

\subsection{17. Relation \#15: Higgs Coupling λ\_H = √17/32}\label{relation-15-higgs-coupling-ux3bb_h-1732}

\textbf{Statement}: The Higgs quartic coupling is determined by G₂ holonomy parameters.

\textbf{Classification}: TOPOLOGICAL

\subsubsection{Proof}\label{proof-19}

\emph{Formula}:
\[\lambda_H = \frac{\sqrt{\dim(G_2) + N_{gen}}}{\det(g)_{den}} = \frac{\sqrt{14 + 3}}{32} = \frac{\sqrt{17}}{32} = 0.12885\]

\emph{Components}:
- dim(G₂) + N\_gen = 17: holonomy dimension plus generation count
- det(g)\_den = 32: metric determinant denominator (= b₂ + dim(G₂) − N\_gen)

\emph{Physical interpretation}:
The Higgs self-coupling combines the holonomy group dimension with the generation count, normalized by the metric determinant scale. The square root reflects the quadratic nature of the Higgs potential.

\emph{Numerical value}: λ\_H = √17/32 = 0.128847\ldots{}

\emph{Experimental comparison}:

{\def\LTcaptype{none} % do not increment counter
\begin{longtable}[]{@{}ll@{}}
\toprule\noalign{}
Quantity & Value \\
\midrule\noalign{}
\endhead
\bottomrule\noalign{}
\endlastfoot
Experimental (PDG 2024) & 0.129 ± 0.001 \\
GIFT prediction & √17/32 = 0.12885 \\
Deviation & 0.12\% \\
\end{longtable}
}

\subsubsection{TCS Alternative}\label{tcs-alternative}

A purely rational TCS formula is also available:

\[\lambda_H^{TCS} = \frac{b_2(M_1)}{b_3 + b_2(M_2)} = \frac{11}{87} = 0.12644\]

where b₂(M₁) = 11 and b₂(M₂) = 10 are the Betti numbers of the two TCS building blocks. This formula is lower-complexity (cost 3 vs cost 6) and connects directly to the twisted connected sum decomposition. Against a running coupling value λ\_H(μ\_high) ≈ 0.126, it achieves 0.35\% deviation, though against the PDG 2024 pole value (0.129) it gives 1.98\%.

\textbf{Status}: TOPOLOGICAL ∎

\begin{center}\rule{0.5\linewidth}{0.5pt}\end{center}

\subsection{17b. Boson Mass Ratios}\label{b.-boson-mass-ratios}

\textbf{Statement}: The ratios of electroweak boson masses have topological origins.

\textbf{Classification}: VERIFIED

\subsubsection{Relation: m\_W/m\_Z = 37/42}\label{relation-m_wm_z-3742}

\emph{Formula}:
\[\frac{m_W}{m_Z} = \frac{2b_2 - \text{Weyl}}{2b_2} = \frac{42 - 5}{42} = \frac{37}{42}\]

\emph{Physical interpretation}:
- 2b₂ = 42 is the structural constant (= p₂ × b₂)
- Weyl = 5 is the triple identity factor
- The ratio involves (structural\_const − Weyl) / structural\_const

\textbf{Note}: The true Euler characteristic χ(K₇) = 0 for odd-dimensional manifolds. The constant 42 = 2b₂ is a distinct topological invariant.

\emph{Numerical value}: m\_W/m\_Z = 0.8810

\emph{Experimental comparison}:

{\def\LTcaptype{none} % do not increment counter
\begin{longtable}[]{@{}ll@{}}
\toprule\noalign{}
Quantity & Value \\
\midrule\noalign{}
\endhead
\bottomrule\noalign{}
\endlastfoot
Experimental & 0.8815 ± 0.0002 \\
GIFT prediction & 0.8810 \\
Deviation & \textbf{0.06\%} \\
\end{longtable}
}

\subsubsection{Relation: m\_H/m\_t = 56/77}\label{relation-m_hm_t-5677}

\emph{Formula}:
\[\frac{m_H}{m_t} = \frac{fund(E_7)}{b_3} = \frac{56}{77} = \frac{8}{11}\]

\emph{Numerical value}: m\_H/m\_t = 0.7273

{\def\LTcaptype{none} % do not increment counter
\begin{longtable}[]{@{}ll@{}}
\toprule\noalign{}
Quantity & Value \\
\midrule\noalign{}
\endhead
\bottomrule\noalign{}
\endlastfoot
Experimental & 0.725 ± 0.003 \\
GIFT prediction & 0.7273 \\
Deviation & 0.31\% \\
\end{longtable}
}

\subsubsection{Relation: m\_H/m\_W = 81/52}\label{relation-m_hm_w-8152}

\emph{Formula}:
\[\frac{m_H}{m_W} = \frac{N_{gen} + \dim(E_6)}{\dim(F_4)} = \frac{3 + 78}{52} = \frac{81}{52}\]

\emph{Numerical value}: m\_H/m\_W = 1.5577

{\def\LTcaptype{none} % do not increment counter
\begin{longtable}[]{@{}ll@{}}
\toprule\noalign{}
Quantity & Value \\
\midrule\noalign{}
\endhead
\bottomrule\noalign{}
\endlastfoot
Experimental & 1.558 ± 0.002 \\
GIFT prediction & 1.5577 \\
Deviation & \textbf{0.02\%} \\
\end{longtable}
}

\textbf{Status}: VERIFIED ∎

\begin{center}\rule{0.5\linewidth}{0.5pt}\end{center}

\subsection{18. Relation \#16: Dark Energy Density Ω\_DE = ln(2)×98/99}\label{relation-16-dark-energy-density-ux3c9_de-ln29899}

\textbf{Statement}: The dark energy density fraction emerges from information-theoretic structure of K₇ moduli.

\textbf{Classification}: TOPOLOGICAL

\subsubsection{Proof}\label{proof-20}

\emph{Formula}:
\[\Omega_{DE} = \ln(2) \times \frac{b_2 + b_3}{H^*} = \ln(2) \times \frac{98}{99} = 0.6861\]

\emph{Components}:
- ln(2): binary entropy factor
- b₂ + b₃ = 98: total non-trivial cohomological content
- H* = 99 = b₂ + b₃ + 1: full cohomological dimension

\emph{Physical interpretation}:
The dark energy fraction combines a binary information factor with the ratio of non-trivial to total cohomological content. The ln(2) reflects the fundamental binary structure of the moduli space partition.

\emph{Numerical value}: Ω\_DE = ln(2) × 98/99 = 0.68615\ldots{}

\emph{Experimental comparison}:

{\def\LTcaptype{none} % do not increment counter
\begin{longtable}[]{@{}ll@{}}
\toprule\noalign{}
Quantity & Value \\
\midrule\noalign{}
\endhead
\bottomrule\noalign{}
\endlastfoot
Experimental (Planck PR4) & 0.6847 ± 0.005 \\
GIFT prediction & ln(2)×98/99 = 0.6861 \\
Deviation & 0.21\% \\
\end{longtable}
}

\subsubsection{TCS Alternative}\label{tcs-alternative-1}

A purely rational TCS formula is also available:

\[\Omega_{DE}^{TCS} = 1 - \frac{\chi(K3)}{b_3} = 1 - \frac{24}{77} = \frac{53}{77} = 0.68831\]

where χ(K3) = 24 is the Euler characteristic of the K3 fiber. This formula is lower-complexity (rational, cost 2 vs cost 6), and the interpretation is transparent: Ω\_matter = χ(K3)/b₃ counts the K3-fiber fraction of the 3-cycle moduli, with the rest being vacuum energy. Against Planck 2018 (Ω\_Λ = 0.6889), it achieves 0.085\% deviation, though against Planck PR4 (0.6847) it gives 0.53\%.

\textbf{Status}: TOPOLOGICAL ∎

\begin{center}\rule{0.5\linewidth}{0.5pt}\end{center}

\subsection{19. Relation \#17: Spectral Index n\_s}\label{relation-17-spectral-index-n_s}

\textbf{Statement}: The primordial scalar spectral index.

\textbf{Classification}: VERIFIED

\subsubsection{Proof}\label{proof-21}

\emph{Formula}:
\[n_s = \frac{\zeta(D_{bulk})}{\zeta(\text{Weyl})} = \frac{\zeta(11)}{\zeta(5)} = 0.9649\]

\emph{Components}:
- ζ(11): From 11D bulk spacetime
- ζ(5): From Weyl factor

\emph{Experimental comparison}:

{\def\LTcaptype{none} % do not increment counter
\begin{longtable}[]{@{}ll@{}}
\toprule\noalign{}
Quantity & Value \\
\midrule\noalign{}
\endhead
\bottomrule\noalign{}
\endlastfoot
Experimental (Planck 2020) & 0.9649 ± 0.0042 \\
GIFT prediction & 0.9649 \\
Deviation & 0.004\% \\
\end{longtable}
}

\textbf{Status}: VERIFIED ∎

\begin{center}\rule{0.5\linewidth}{0.5pt}\end{center}

\subsection{19b. Relation \#17c: Dark Matter to Baryon Ratio Ω\_DM/Ω\_b = 43/8}\label{b.-relation-17c-dark-matter-to-baryon-ratio-ux3c9_dmux3c9_b-438}

\textbf{Statement}: The dark matter to baryon density ratio.

\textbf{Classification}: TOPOLOGICAL

\subsubsection{Proof}\label{proof-22}

\emph{Formula}:
\[\frac{\Omega_{DM}}{\Omega_b} = \frac{b_0 + 42}{\text{rank}(E_8)} = \frac{1 + 42}{8} = \frac{43}{8} = 5.375\]

\emph{Components}:
- 42 = p₂ × N\_gen × dim(K₇): The same constant appearing in m\_b/m\_t = 1/42
- rank(E₈) = 8: Cartan subalgebra dimension

\emph{Experimental comparison}:

{\def\LTcaptype{none} % do not increment counter
\begin{longtable}[]{@{}ll@{}}
\toprule\noalign{}
Quantity & Value \\
\midrule\noalign{}
\endhead
\bottomrule\noalign{}
\endlastfoot
Experimental (Planck 2020) & 5.375 ± 0.05 \\
GIFT prediction & 5.375 \\
Deviation & 0.00\% \\
\end{longtable}
}

\textbf{Status}: TOPOLOGICAL ∎

\begin{center}\rule{0.5\linewidth}{0.5pt}\end{center}

\subsection{19c. Relation \#17d: Reduced Hubble Parameter h = 167/248}\label{c.-relation-17d-reduced-hubble-parameter-h-167248}

\textbf{Statement}: The reduced Hubble parameter H₀ = 100h km/s/Mpc.

\textbf{Classification}: TOPOLOGICAL

\subsubsection{Proof}\label{proof-23}

\emph{Formula}:
\[h = \frac{|PSL_2(7)| - b_0}{\dim(E_8)} = \frac{168 - 1}{248} = \frac{167}{248} = 0.6734...\]

\emph{Experimental comparison}:

{\def\LTcaptype{none} % do not increment counter
\begin{longtable}[]{@{}ll@{}}
\toprule\noalign{}
Quantity & Value \\
\midrule\noalign{}
\endhead
\bottomrule\noalign{}
\endlastfoot
Experimental (Planck 2020) & 0.674 ± 0.005 \\
GIFT prediction & 0.6734 \\
Deviation & 0.09\% \\
\end{longtable}
}

\textbf{Status}: TOPOLOGICAL ∎

\begin{center}\rule{0.5\linewidth}{0.5pt}\end{center}

\subsection{19d. Relation \#17e: Baryon Fraction Ω\_b/Ω\_m = 5/32}\label{d.-relation-17e-baryon-fraction-ux3c9_bux3c9_m-532}

\textbf{Statement}: The baryon fraction of total matter.

\textbf{Classification}: TOPOLOGICAL

\subsubsection{Proof}\label{proof-24}

\emph{Formula}:
\[\frac{\Omega_b}{\Omega_m} = \frac{\text{Weyl}}{\det(g)_{den}} = \frac{5}{32} = 0.15625\]

\emph{Experimental comparison}:

{\def\LTcaptype{none} % do not increment counter
\begin{longtable}[]{@{}ll@{}}
\toprule\noalign{}
Quantity & Value \\
\midrule\noalign{}
\endhead
\bottomrule\noalign{}
\endlastfoot
Experimental (Planck 2020) & 0.156 ± 0.003 \\
GIFT prediction & 0.15625 \\
Deviation & 0.16\% \\
\end{longtable}
}

\textbf{Status}: TOPOLOGICAL ∎

\begin{center}\rule{0.5\linewidth}{0.5pt}\end{center}

\subsection{19e. Relation \#17f: Amplitude of Fluctuations σ₈ = 17/21}\label{e.-relation-17f-amplitude-of-fluctuations-ux3c3ux2088-1721}

\textbf{Statement}: The amplitude of matter fluctuations at 8 h⁻¹ Mpc.

\textbf{Classification}: TOPOLOGICAL

\subsubsection{Proof}\label{proof-25}

\emph{Formula}:
\[\sigma_8 = \frac{p_2 + \det(g)_{den}}{42} = \frac{2 + 32}{42} = \frac{34}{42} = \frac{17}{21} = 0.8095...\]

\emph{Experimental comparison}:

{\def\LTcaptype{none} % do not increment counter
\begin{longtable}[]{@{}ll@{}}
\toprule\noalign{}
Quantity & Value \\
\midrule\noalign{}
\endhead
\bottomrule\noalign{}
\endlastfoot
Experimental (Planck 2020) & 0.811 ± 0.006 \\
GIFT prediction & 0.8095 \\
Deviation & 0.18\% \\
\end{longtable}
}

\textbf{Status}: TOPOLOGICAL ∎

\begin{center}\rule{0.5\linewidth}{0.5pt}\end{center}

\subsection{19f. Relation \#17g: Primordial Helium Fraction Y\_p = 15/61}\label{f.-relation-17g-primordial-helium-fraction-y_p-1561}

\textbf{Statement}: The primordial helium mass fraction from Big Bang nucleosynthesis.

\textbf{Classification}: TOPOLOGICAL

\subsubsection{Proof}\label{proof-26}

\emph{Formula}:
\[Y_p = \frac{b_0 + \dim(G_2)}{\kappa_T^{-1}} = \frac{1 + 14}{61} = \frac{15}{61} = 0.2459...\]

\emph{Experimental comparison}:

{\def\LTcaptype{none} % do not increment counter
\begin{longtable}[]{@{}ll@{}}
\toprule\noalign{}
Quantity & Value \\
\midrule\noalign{}
\endhead
\bottomrule\noalign{}
\endlastfoot
Experimental & 0.245 ± 0.003 \\
GIFT prediction & 0.2459 \\
Deviation & 0.37\% \\
\end{longtable}
}

\textbf{Status}: TOPOLOGICAL ∎

\begin{center}\rule{0.5\linewidth}{0.5pt}\end{center}

\subsection{20. Relation \#17b: Matter Density Ω\_m}\label{relation-17b-matter-density-ux3c9_m}

\textbf{Statement}: The matter density fraction derives from dark energy via √Weyl.

\textbf{Classification}: DERIVED (from Weyl triple identity + Ω\_DE)

\subsubsection{Proof}\label{proof-27}

\emph{Step 1: Establish √Weyl as structural}

From the Weyl Triple Identity (S1, Section 2.3):
\[\text{Weyl} = \frac{\dim(G_2) + 1}{N_{gen}} = \frac{b_2}{N_{gen}} - p_2 = \dim(G_2) - \text{rank}(E_8) - 1 = 5\]

Therefore √Weyl = √5 is a derived quantity.

\emph{Step 2: Matter-dark energy ratio}

The cosmological density ratio:
\[\frac{\Omega_{DE}}{\Omega_m} = \sqrt{\text{Weyl}} = \sqrt{5}\]

\emph{Step 3: Compute Ω\_m}

Using Ω\_DE = ln(2) × (b₂ + b₃)/H* = 0.6861 (Relation \#16):
\[\Omega_m = \frac{\Omega_{DE}}{\sqrt{\text{Weyl}}} = \frac{\ln(2) \times 98/99}{\sqrt{5}} = \frac{0.6861}{2.236} = 0.3068\]

\emph{Step 4: Verify closure}

\[\Omega_{total} = \Omega_{DE} + \Omega_m = 0.6861 + 0.3068 = 0.9929 \approx 1\]

Consistent with flat universe (Ω\_total = 1).

\emph{Experimental comparison}:

{\def\LTcaptype{none} % do not increment counter
\begin{longtable}[]{@{}ll@{}}
\toprule\noalign{}
Quantity & Value \\
\midrule\noalign{}
\endhead
\bottomrule\noalign{}
\endlastfoot
Experimental (Planck 2020) & 0.3153 ± 0.007 \\
GIFT prediction & 0.3068 \\
Deviation & 2.7\% \\
\end{longtable}
}

\subsubsection{Interpretation}\label{interpretation}

The √5 ratio between dark energy and matter densities emerges from the same structural constant (Weyl = 5) that determines:
- det(g) = 65/32 (metric determinant)
- \textbar W(E₈)\textbar{} factorization (group theory)
- N\_gen³ coefficient in \textbar W(E₈)\textbar{} (topology)

\textbf{Status}: DERIVED (structural, 2.7\% deviation) ∎

\begin{center}\rule{0.5\linewidth}{0.5pt}\end{center}

\subsection{21. Relation \#18: Fine Structure Constant α⁻¹}\label{relation-18-fine-structure-constant-ux3b1uxb9}

\textbf{Statement}: The inverse fine structure constant.

\textbf{Classification}: TOPOLOGICAL

\subsubsection{Proof}\label{proof-28}

\emph{Formula}:
\[\alpha^{-1}(M_Z) = \frac{\dim(E_8) + \text{rank}(E_8)}{2} + \frac{H^*}{D_{bulk}} + \det(g) \cdot \kappa_T\]
\[= 128 + 9 + \frac{65}{32} \times \frac{1}{61} = 137.033\]

\emph{Components}:
- 128 = (248 + 8)/2: Algebraic
- 9 = 99/11: Bulk impedance
- 65/1952: Torsional correction

\emph{Experimental comparison}:

{\def\LTcaptype{none} % do not increment counter
\begin{longtable}[]{@{}ll@{}}
\toprule\noalign{}
Quantity & Value \\
\midrule\noalign{}
\endhead
\bottomrule\noalign{}
\endlastfoot
Experimental & 137.035999 \\
GIFT prediction & 137.033 \\
Deviation & 0.002\% \\
\end{longtable}
}

\textbf{Status}: TOPOLOGICAL ∎

\begin{center}\rule{0.5\linewidth}{0.5pt}\end{center}

\section{Part VIII: Summary Table}\label{part-viii-summary-table}

\subsection{21. The 33 Type I Dimensionless Relations (+ 2 Type IV BH Appendix)}\label{the-33-type-i-dimensionless-relations-2-type-iv-bh-appendix}

\textbf{Note}: All predictions use only topological invariants (b\textsubscript{2}, b\textsubscript{3}, dim(G\textsubscript{2}), etc.). None depend on the realized torsion value T. Relation \#19 (Ω\_m) is DERIVED from Ω\_DE via the Weyl triple identity.

{\def\LTcaptype{none} % do not increment counter
\begin{longtable}[]{@{}
  >{\raggedright\arraybackslash}p{(\linewidth - 12\tabcolsep) * \real{0.0612}}
  >{\raggedright\arraybackslash}p{(\linewidth - 12\tabcolsep) * \real{0.2041}}
  >{\raggedright\arraybackslash}p{(\linewidth - 12\tabcolsep) * \real{0.1837}}
  >{\raggedright\arraybackslash}p{(\linewidth - 12\tabcolsep) * \real{0.1429}}
  >{\raggedright\arraybackslash}p{(\linewidth - 12\tabcolsep) * \real{0.1224}}
  >{\raggedright\arraybackslash}p{(\linewidth - 12\tabcolsep) * \real{0.1224}}
  >{\raggedright\arraybackslash}p{(\linewidth - 12\tabcolsep) * \real{0.1633}}@{}}
\toprule\noalign{}
\begin{minipage}[b]{\linewidth}\raggedright
\#
\end{minipage} & \begin{minipage}[b]{\linewidth}\raggedright
Relation
\end{minipage} & \begin{minipage}[b]{\linewidth}\raggedright
Formula
\end{minipage} & \begin{minipage}[b]{\linewidth}\raggedright
Value
\end{minipage} & \begin{minipage}[b]{\linewidth}\raggedright
Exp.
\end{minipage} & \begin{minipage}[b]{\linewidth}\raggedright
Dev.
\end{minipage} & \begin{minipage}[b]{\linewidth}\raggedright
Status
\end{minipage} \\
\midrule\noalign{}
\endhead
\bottomrule\noalign{}
\endlastfoot
1 & N\_gen & Atiyah-Singer & 3 & 3 & exact & VERIFIED \\
2 & τ & 496×21/(27×99) & 3472/891 & - & - & VERIFIED \\
3 & κ\_T & 1/(77-14-2) & 1/61 & - & - & STRUCTURAL* \\
4 & det(g) & 5×13/32 & 65/32 & - & - & METRIC NORM. \\
5 & sin²θ\_W & 21/91 & 3/13 & 0.23122 & 0.195\% & VERIFIED \\
6 & α\_s & √2/12 & 0.11785 & 0.1180 & 0.127\% & TOPOLOGICAL \\
7 & Q\_Koide & 14/21 & 2/3 & 0.666661 & 0.0009\% & VERIFIED \\
8 & m\_τ/m\_e & 7+2480+990 & 3477 & 3477.15 & 0.0043\% & VERIFIED \\
9 & m\_μ/m\_e & 27\^{}φ & 207.01 & 206.768 & 0.118\% & TOPOLOGICAL \\
10 & m\_s/m\_d & 4×5 & 20 & 20.0 & 0.00\% & VERIFIED \\
11 & δ\_CP & 7×14+99 & 197° & 177°±20°* & 11.3\% (1σ) & VERIFIED \\
12 & θ₁₃ & π/21 (TCS: 94/11) & 8.57° & 8.54° & 0.37\% & TOPOLOGICAL \\
13 & θ₂₃ & arcsin((b₃-p₂)/H*) & 49.25° & 49.3° & 0.10\% & TOPOLOGICAL \\
14 & θ₁₂ & arctan(\ldots) & 33.40° & 33.41° & 0.030\% & TOPOLOGICAL \\
15 & λ\_H & √(dim(G₂)+N\_gen)/det\_den & √17/32 & 0.129 & 0.12\% & TOPOLOGICAL \\
16 & Ω\_DE & ln(2)×(b₂+b₃)/H* & 0.6861 & 0.6847 & 0.21\% & TOPOLOGICAL \\
17 & n\_s & ζ(11)/ζ(5) & 0.9649 & 0.9649 & 0.004\% & VERIFIED \\
18 & α⁻¹ & 128+9+corr & 137.033 & 137.036 & 0.002\% & TOPOLOGICAL \\
19 & Ω\_m & Ω\_DE/√Weyl & 0.3068 & 0.3153 & 2.7\% & DERIVED \\
\end{longtable}
}

*NuFIT 6.0 (Oct 2024); ±20° uncertainty. The 197° prediction matched NuFIT 5.2 (2022) exactly. See §4.4.1 of main paper.

*κ\_T is a structural parameter (capacity), not a physical prediction. It does not appear in other formulas.

\begin{center}\rule{0.5\linewidth}{0.5pt}\end{center}

\subsection{22. Deviation Statistics}\label{deviation-statistics}

{\def\LTcaptype{none} % do not increment counter
\begin{longtable}[]{@{}lll@{}}
\toprule\noalign{}
Range & Count & Percentage \\
\midrule\noalign{}
\endhead
\bottomrule\noalign{}
\endlastfoot
0.00\% (exact) & 4 & 22\% \\
\textless0.01\% & 3 & 17\% \\
0.01-0.1\% & 4 & 22\% \\
0.1-0.5\% & 7 & 39\% \\
\end{longtable}
}

\textbf{Mean deviation}: 0.73\% (PDG 2024, NuFIT 6.0; 33 Type I observables with experimental comparison, canonical δ\_CP=197°). For the full 92-observable dataset (66 with experimental comparison), see Supplement S3.

\begin{center}\rule{0.5\linewidth}{0.5pt}\end{center}

\subsection{23. Statistical Uniqueness of (b₂=21, b₃=77)}\label{statistical-uniqueness-of-bux208221-bux208377}

A critical question for any unified framework is whether the specific topological parameters represent overfitting. We conducted comprehensive Monte Carlo validation to address this concern.

\subsubsection{Methodology}\label{methodology}

\begin{itemize}
\item
  \textbf{Betti variations}: 100,000 random (b₂, b₃) configurations
\item
  \textbf{Gauge group comparison}: E₈×E₈, E₇×E₇, E₆×E₆, SO(32), SU(5)×SU(5), etc.
\item
  \textbf{Holonomy comparison}: G₂, Spin(7), SU(3), SU(4)
\item
  \textbf{Full combinatorial}: 91,896 configurations varying all parameters
\item
  \textbf{Local sensitivity}: ±10 grid around (b₂=21, b₃=77)
\end{itemize}

\subsubsection{Results}\label{results}

{\def\LTcaptype{none} % do not increment counter
\begin{longtable}[]{@{}ll@{}}
\toprule\noalign{}
Metric & Value \\
\midrule\noalign{}
\endhead
\bottomrule\noalign{}
\endlastfoot
Total configurations tested & \textbf{192,349} \\
Configurations better than GIFT & \textbf{0} \\
GIFT mean deviation & 0.73\% (33 observables) \\
Alternative mean deviation & 32.9\% \\
P-value & \textbf{\textless{} 5 × 10⁻⁶} \\
Significance & \textbf{\textgreater{} 4.5σ} \\
\end{longtable}
}

\subsubsection{Gauge Group Ranking}\label{gauge-group-ranking}

{\def\LTcaptype{none} % do not increment counter
\begin{longtable}[]{@{}lll@{}}
\toprule\noalign{}
Rank & Group & Mean Dev. \\
\midrule\noalign{}
\endhead
\bottomrule\noalign{}
\endlastfoot
\textbf{1} & \textbf{E₈×E₈} & \textbf{0.73\%} \\
2 & E₇×E₈ & 8.80\% \\
3 & E₆×E₈ & 15.50\% \\
\end{longtable}
}

E₈ x E₈ achieves approximately 10x better agreement than all tested alternatives.

\subsubsection{Holonomy Ranking}\label{holonomy-ranking}

{\def\LTcaptype{none} % do not increment counter
\begin{longtable}[]{@{}lll@{}}
\toprule\noalign{}
Rank & Holonomy & Mean Dev. \\
\midrule\noalign{}
\endhead
\bottomrule\noalign{}
\endlastfoot
\textbf{1} & \textbf{G₂} & \textbf{0.73\%} \\
2 & SU(4) & 1.46\% \\
3 & SU(3) & 4.43\% \\
\end{longtable}
}

G₂ achieves approximately 5x better agreement than Calabi-Yau (SU(3)).

\subsubsection{Interpretation}\label{interpretation-1}

Among 192,349 tested alternatives, the configuration (b₂=21, b₃=77) with E₈ x E₈ gauge group and G₂ holonomy achieves the lowest mean deviation. Zero alternatives outperform it.

\textbf{Literature uniqueness}: Analysis of the NK-certified metric ([B], §6) compared against all \textasciitilde65 known compact G₂ manifolds from the literature (Kovalev, CHNP, CGN, Joyce, Nordström). The point (b₂=21, b₃=77) is UNIQUE: no known diffeomorphic example exists. The nearest neighbor in (b₂,b₃) space is the CHNP grid point (9,9), at distance 7.6. The topological pair is realized by the Joyce-Karigiannis Z₂³ orbifold route (S1 §8.4), extending the known catalogue; an explicit closed-form neck ansatz is established in [D].

Complete methodology: \href{../../docs/STATISTICAL_EVIDENCE.md}{docs/STATISTICAL\_EVIDENCE.md}

\begin{center}\rule{0.5\linewidth}{0.5pt}\end{center}

\section{Part IX: Observable Catalog}\label{part-ix-observable-catalog}

\subsection{24. Structural Redundancy and Expression Counts}\label{structural-redundancy-and-expression-counts}

Each prediction admits multiple algebraically independent expressions that reduce to the same fraction. This multiplicity provides a measure of structural robustness: quantities arising from many paths through the topological invariants are less likely to represent numerical coincidence.

\subsubsection{24.1 Classification Scheme}\label{classification-scheme}

{\def\LTcaptype{none} % do not increment counter
\begin{longtable}[]{@{}
  >{\raggedright\arraybackslash}p{(\linewidth - 4\tabcolsep) * \real{0.3556}}
  >{\raggedright\arraybackslash}p{(\linewidth - 4\tabcolsep) * \real{0.2889}}
  >{\raggedright\arraybackslash}p{(\linewidth - 4\tabcolsep) * \real{0.3556}}@{}}
\toprule\noalign{}
\begin{minipage}[b]{\linewidth}\raggedright
Classification
\end{minipage} & \begin{minipage}[b]{\linewidth}\raggedright
Expressions
\end{minipage} & \begin{minipage}[b]{\linewidth}\raggedright
Interpretation
\end{minipage} \\
\midrule\noalign{}
\endhead
\bottomrule\noalign{}
\endlastfoot
\textbf{CANONICAL} & ≥20 & Maximally over-determined; emerges from algebraic web \\
\textbf{ROBUST} & 10-19 & Highly constrained; multiple independent derivations \\
\textbf{SUPPORTED} & 5-9 & Structural redundancy \\
\textbf{DERIVED} & 2-4 & Dual derivation minimum \\
\textbf{SINGULAR} & 1 & Unique path (possible coincidence) \\
\end{longtable}
}

\subsubsection{24.2 Core 18 Type I Predictions with Expression Counts (subset of 33 certified)}\label{core-18-type-i-predictions-with-expression-counts-subset-of-33-certified}

{\def\LTcaptype{none} % do not increment counter
\begin{longtable}[]{@{}
  >{\raggedright\arraybackslash}p{(\linewidth - 14\tabcolsep) * \real{0.0526}}
  >{\raggedright\arraybackslash}p{(\linewidth - 14\tabcolsep) * \real{0.2105}}
  >{\raggedright\arraybackslash}p{(\linewidth - 14\tabcolsep) * \real{0.1579}}
  >{\raggedright\arraybackslash}p{(\linewidth - 14\tabcolsep) * \real{0.1228}}
  >{\raggedright\arraybackslash}p{(\linewidth - 14\tabcolsep) * \real{0.1053}}
  >{\raggedright\arraybackslash}p{(\linewidth - 14\tabcolsep) * \real{0.1053}}
  >{\raggedright\arraybackslash}p{(\linewidth - 14\tabcolsep) * \real{0.1228}}
  >{\raggedright\arraybackslash}p{(\linewidth - 14\tabcolsep) * \real{0.1228}}@{}}
\toprule\noalign{}
\begin{minipage}[b]{\linewidth}\raggedright
\#
\end{minipage} & \begin{minipage}[b]{\linewidth}\raggedright
Observable
\end{minipage} & \begin{minipage}[b]{\linewidth}\raggedright
Formula
\end{minipage} & \begin{minipage}[b]{\linewidth}\raggedright
Value
\end{minipage} & \begin{minipage}[b]{\linewidth}\raggedright
Exp.
\end{minipage} & \begin{minipage}[b]{\linewidth}\raggedright
Dev.
\end{minipage} & \begin{minipage}[b]{\linewidth}\raggedright
Expr.
\end{minipage} & \begin{minipage}[b]{\linewidth}\raggedright
Class
\end{minipage} \\
\midrule\noalign{}
\endhead
\bottomrule\noalign{}
\endlastfoot
1 & N\_gen & Atiyah-Singer & 3 & 3 & 0.00\% & 24+ & CANONICAL \\
2 & sin²θ\_W & b₂/(b₃+dim\_G₂) & 3/13 & 0.2312 & 0.20\% & 14 & ROBUST \\
3 & α\_s(M\_Z) & √2/12 & 0.11785 & 0.11800 & 0.126\% & 9 & SUPPORTED \\
4 & λ\_H & √17/32 & 0.1288 & 0.129 & 0.12\% & 4 & DERIVED \\
5 & α⁻¹ & 128+9+corr & 137.033 & 137.036 & 0.002\% & 3 & DERIVED \\
6 & Q\_Koide & dim\_G₂/b₂ & 2/3 & 0.6667 & 0.001\% & 20 & CANONICAL \\
7 & m\_τ/m\_e & 7+10×248+10×99 & 3477 & 3477.2 & 0.004\% & 3 & DERIVED \\
8 & m\_μ/m\_e & 27\^{}φ & 207.01 & 206.77 & 0.12\% & 2 & DERIVED \\
9 & m\_s/m\_d & p₂²×Weyl & 20 & 20.0 & 0.00\% & 14 & ROBUST \\
10 & m\_b/m\_t & 1/(2b₂) & 1/42 & 0.024 & 0.79\% & 21 & CANONICAL \\
11 & m\_u/m\_d & (1+dim\_E₆)/PSL₂₇ & 79/168 & 0.47 & 0.05\% & 1 & SINGULAR \\
12 & δ\_CP & 7×14+99 & 197° & 177°±20°* & 11.3\% (1σ) & 3 & VERIFIED \\
13 & θ₁₃ & π/b₂ & 8.57° & 8.54° & 0.37\% & 3 & DERIVED \\
14 & θ₂₃ & arcsin((b₃-p₂)/H*) & 49.25° & 49.3° & 0.10\% & 2 & DERIVED \\
15 & θ₁₂ & arctan(√(δ/γ)) & 33.40° & 33.41° & 0.03\% & 2 & DERIVED \\
16 & Ω\_DE & ln(2)×(b₂+b₃)/H* & 0.6861 & 0.6847 & 0.21\% & 2 & DERIVED \\
17 & n\_s & ζ(11)/ζ(5) & 0.9649 & 0.9649 & 0.004\% & 2 & DERIVED \\
18 & det(g) & 65/32 & 2.0313 & - & - & 8 & SUPPORTED \\
\end{longtable}
}

\textbf{Distribution}: 4 CANONICAL (22\%), 4 ROBUST (22\%), 2 SUPPORTED (11\%), 7 DERIVED (39\%), 1 SINGULAR (6\%).

\subsubsection{24.3 Extended Certified Predictions (\#19--\#33, 15 entries)}\label{extended-certified-predictions-1933-15-entries}

{\def\LTcaptype{none} % do not increment counter
\begin{longtable}[]{@{}
  >{\raggedright\arraybackslash}p{(\linewidth - 14\tabcolsep) * \real{0.0526}}
  >{\raggedright\arraybackslash}p{(\linewidth - 14\tabcolsep) * \real{0.2105}}
  >{\raggedright\arraybackslash}p{(\linewidth - 14\tabcolsep) * \real{0.1579}}
  >{\raggedright\arraybackslash}p{(\linewidth - 14\tabcolsep) * \real{0.1228}}
  >{\raggedright\arraybackslash}p{(\linewidth - 14\tabcolsep) * \real{0.1053}}
  >{\raggedright\arraybackslash}p{(\linewidth - 14\tabcolsep) * \real{0.1053}}
  >{\raggedright\arraybackslash}p{(\linewidth - 14\tabcolsep) * \real{0.1228}}
  >{\raggedright\arraybackslash}p{(\linewidth - 14\tabcolsep) * \real{0.1228}}@{}}
\toprule\noalign{}
\begin{minipage}[b]{\linewidth}\raggedright
\#
\end{minipage} & \begin{minipage}[b]{\linewidth}\raggedright
Observable
\end{minipage} & \begin{minipage}[b]{\linewidth}\raggedright
Formula
\end{minipage} & \begin{minipage}[b]{\linewidth}\raggedright
Value
\end{minipage} & \begin{minipage}[b]{\linewidth}\raggedright
Exp.
\end{minipage} & \begin{minipage}[b]{\linewidth}\raggedright
Dev.
\end{minipage} & \begin{minipage}[b]{\linewidth}\raggedright
Expr.
\end{minipage} & \begin{minipage}[b]{\linewidth}\raggedright
Class
\end{minipage} \\
\midrule\noalign{}
\endhead
\bottomrule\noalign{}
\endlastfoot
19 & sin²θ₁₂\^{}PMNS & (1+N\_gen)/α\_sum & 4/13 & 0.307 & 0.23\% & 28 & CANONICAL \\
20 & sin²θ₂₃\^{}PMNS & (D\_bulk−Weyl)/D\_bulk & 6/11 & 0.546 & 0.10\% & 15 & ROBUST \\
21 & sin²θ₁₃\^{}PMNS & D\_bulk/dim\_E₈₂ & 11/496 & 0.022 & 0.81\% & 5 & SUPPORTED \\
22 & sin²θ₁₂\^{}CKM & 7/31 & 0.2258 & 0.225 & 0.36\% & 16 & ROBUST \\
23 & A\_Wolf & (Weyl+dim\_E₆)/H* & 83/99 & 0.836 & 0.29\% & 4 & DERIVED \\
24 & sin²θ₂₃\^{}CKM & dim\_K₇/PSL₂₇ & 1/24 & 0.041 & 1.13\% & 3 & DERIVED \\
25 & m\_H/m\_t & 8/11 & 0.7273 & 0.725 & 0.31\% & 19 & ROBUST \\
26 & m\_H/m\_W & 81/52 & 1.5577 & 1.558 & 0.02\% & 1 & SINGULAR \\
27 & m\_W/m\_Z & (2b₂−Weyl)/(2b₂) = 37/42 & 0.8810 & 0.8815 & \textbf{0.06\%} & 8 & SUPPORTED \\
28 & m\_μ/m\_τ & 5/84 & 0.0595 & 0.0595 & 0.04\% & 9 & SUPPORTED \\
29 & Ω\_DM/Ω\_b & (1+42)/rank\_E₈ & 43/8 & 5.375 & 0.00\% & 6 & SUPPORTED \\
30 & Ω\_b/Ω\_m & Weyl/det(g)\_den & 5/32 & 0.156 & 0.16\% & 7 & SUPPORTED \\
31 & Ω\_Λ/Ω\_m & (det\_g\_den−dim\_K₇)/D\_bulk & 25/11 & 2.27 & 0.12\% & 6 & SUPPORTED \\
32 & h & (PSL₂₇−1)/dim\_E₈ & 167/248 & 0.674 & 0.09\% & 3 & DERIVED \\
33 & σ₈ & (p₂+det\_g\_den)/(2b₂) & 34/42 & 0.811 & 0.18\% & 4 & DERIVED \\
\end{longtable}
}

\subsubsection{24.3b Relations \#34--\#35: BH Remnant Topological Predictions: Type IV (Pinčák et al.~2026 \cite{42})}\label{b-relations-3435-bh-remnant-topological-predictions-type-iv-pinux10duxe1k-et-al.-2026-46}

\emph{These are classified as \textbf{Type IV structural diagnostics} (D22 = M\_res, D20 = N\_QNM) in the 95-observable dataset, not as Type I. They are included here for completeness as they are topologically derived from the same G₂ structural constants.}

{\def\LTcaptype{none} % do not increment counter
\begin{longtable}[]{@{}
  >{\raggedright\arraybackslash}p{(\linewidth - 8\tabcolsep) * \real{0.0769}}
  >{\raggedright\arraybackslash}p{(\linewidth - 8\tabcolsep) * \real{0.3077}}
  >{\raggedright\arraybackslash}p{(\linewidth - 8\tabcolsep) * \real{0.2308}}
  >{\raggedright\arraybackslash}p{(\linewidth - 8\tabcolsep) * \real{0.1795}}
  >{\raggedright\arraybackslash}p{(\linewidth - 8\tabcolsep) * \real{0.2051}}@{}}
\toprule\noalign{}
\begin{minipage}[b]{\linewidth}\raggedright
\#
\end{minipage} & \begin{minipage}[b]{\linewidth}\raggedright
Observable
\end{minipage} & \begin{minipage}[b]{\linewidth}\raggedright
Formula
\end{minipage} & \begin{minipage}[b]{\linewidth}\raggedright
Value
\end{minipage} & \begin{minipage}[b]{\linewidth}\raggedright
Status
\end{minipage} \\
\midrule\noalign{}
\endhead
\bottomrule\noalign{}
\endlastfoot
34 & BH remnant mass M\_res & v\_EW²/M\_Pl (torsion VEV: τ₀ = v\_EW) & \textasciitilde125 GeV² / M\_Pl & Type IV; topological, not yet Lean-certified \\
35 & QNM families N\_QNM & b₃ + Weyl×dim(K₇) + p₂ (G₂ torsion modes) & 98 & Type IV; topological, not yet Lean-certified \\
\end{longtable}
}

These two predictions are topologically derived from G₂ torsion structure but have not yet been formally certified in Lean. M\_res involves a physical scale identification (τ₀ = v\_EW); N\_QNM is a purely structural count. No experimental comparison is currently available.

\subsubsection{24.4 Illustrative Examples of Multiple Expressions}\label{illustrative-examples-of-multiple-expressions}

\textbf{sin²θ\_W = 3/13} (14 independent expressions):

{\def\LTcaptype{none} % do not increment counter
\begin{longtable}[]{@{}lll@{}}
\toprule\noalign{}
\# & Expression & Evaluation \\
\midrule\noalign{}
\endhead
\bottomrule\noalign{}
\endlastfoot
1 & N\_gen / α\_sum & 3/13 \\
2 & N\_gen / (p₂ + D\_bulk) & 3/(2+11) = 3/13 \\
3 & b₂ / (b₃ + dim\_G₂) & 21/91 = 3/13 \\
4 & dim(J₃O) / (dim\_F₄ + det\_g\_num) & 27/117 = 3/13 \\
5 & (b₀ + dim\_G₂) / det\_g\_num & 15/65 = 3/13 \\
6 & (p₂ + b₀) / α\_sum & 3/13 \\
7 & dim\_K₇ / (b₂ + dim\_K₇ + dim\_G₂) & 7/42 ≠ 3/13 ✗ \\
\end{longtable}
}

(Expression 7 illustrates that not all combinations work; only those reducing to 3/13 are valid.)

\textbf{Q\_Koide = 2/3} (20 independent expressions):

{\def\LTcaptype{none} % do not increment counter
\begin{longtable}[]{@{}lll@{}}
\toprule\noalign{}
\# & Expression & Evaluation \\
\midrule\noalign{}
\endhead
\bottomrule\noalign{}
\endlastfoot
1 & p₂ / N\_gen & 2/3 \\
2 & dim\_G₂ / b₂ & 14/21 = 2/3 \\
3 & dim\_F₄ / dim\_E₆ & 52/78 = 2/3 \\
4 & rank\_E₈ / (Weyl + dim\_K₇) & 8/12 = 2/3 \\
5 & (dim\_G₂ − rank\_E₈) / (rank\_E₈ + 1) & 6/9 = 2/3 \\
\end{longtable}
}

\textbf{m\_b/m\_t = 1/42} (21 independent expressions):

{\def\LTcaptype{none} % do not increment counter
\begin{longtable}[]{@{}lll@{}}
\toprule\noalign{}
\# & Expression & Evaluation \\
\midrule\noalign{}
\endhead
\bottomrule\noalign{}
\endlastfoot
1 & b₀ / (2b₂) & 1/42 \\
2 & (b₀ + N\_gen) / PSL(2,7) & 4/168 = 1/42 \\
3 & p₂ / (dim\_K₇ + b₃) & 2/84 = 1/42 \\
4 & N\_gen / (dim(J₃O) + H*) & 3/126 = 1/42 \\
5 & dim\_K₇ / (dim\_E₈ + dim(J₃O) + dim\_K₇) & 7/294 = 1/42 \\
\end{longtable}
}

The ratio m\_b/m\_t = 1/42 = 1/(2b₂) illustrates structural redundancy: the bottom-to-top mass hierarchy equals the inverse of the structural constant 2b₂ = p₂ × b₂.

\textbf{Note}: The true Euler characteristic χ(K₇) = 0 for G₂ manifolds (odd-dimensional). The constant 42 is the structural invariant 2b₂.

\subsubsection{24.5 The Algebraic Web}\label{the-algebraic-web}

The topological constants satisfy interconnected identities:

{\def\LTcaptype{none} % do not increment counter
\begin{longtable}[]{@{}lll@{}}
\toprule\noalign{}
Identity & Left side & Right side \\
\midrule\noalign{}
\endhead
\bottomrule\noalign{}
\endlastfoot
Fiber-holonomy & dim(G₂) = 14 & p₂ × dim(K₇) = 2 × 7 \\
Gauge moduli & b₂ = 21 & N\_gen × dim(K₇) = 3 × 7 \\
Matter-holonomy & b₃ + dim(G₂) = 91 & dim(K₇) × α\_sum = 7 × 13 \\
Fano order & PSL(2,7) = 168 & rank(E₈) × b₂ = 8 × 21 \\
Fano order & PSL(2,7) = 168 & N\_gen × fund(E₇) = 3 × 56 \\
Anomaly sum & α\_sum = 13 & rank(E₈) + Weyl = 8 + 5 \\
\end{longtable}
}

These relations form a closed algebraic system. The mod-7 structure (dim(K₇) = 7 divides dim(G₂), b₂, b₃, PSL(2,7)) reflects the Fano plane underlying octonion multiplication.

\subsubsection{24.6 Fibonacci-Lucas Embedding}\label{fibonacci-lucas-embedding}

The GIFT constants embed naturally into the Fibonacci (Fₙ) and Lucas (Lₙ) sequences:

{\def\LTcaptype{none} % do not increment counter
\begin{longtable}[]{@{}llll@{}}
\toprule\noalign{}
n & Fₙ & GIFT Constant & Role \\
\midrule\noalign{}
\endhead
\bottomrule\noalign{}
\endlastfoot
3 & 2 & p₂ & Pontryagin class \\
4 & 3 & N\_gen & Fermion generations \\
5 & 5 & Weyl & Pentagonal symmetry \\
6 & 8 & rank(E₈) & Cartan subalgebra \\
7 & 13 & α²\_B sum & Structure coefficient \\
8 & 21 & b₂ & Second Betti number \\
\end{longtable}
}

This sequence propagates via the recurrence:

\[F_3 + F_4 = F_5 \quad \Rightarrow \quad p_2 + N_{gen} = \text{Weyl}\]

Lucas numbers also appear naturally:

{\def\LTcaptype{none} % do not increment counter
\begin{longtable}[]{@{}lll@{}}
\toprule\noalign{}
Lₙ & Value & GIFT Role \\
\midrule\noalign{}
\endhead
\bottomrule\noalign{}
\endlastfoot
L₄ & 7 & dim(K₇) \\
L₅ & 11 & D\_bulk \\
L₈ & 47 & Scale bridge exponent \\
\end{longtable}
}

The Lucas identity L₈ = F₇ + F₉ = 13 + 34 decomposes as:

\[L_8 = \alpha_{sum}^B + d_{hidden} = 13 + 34 = 47\]

This structure reflects the icosahedral geometry underlying the McKay correspondence E₈ ↔ 2I, where icosahedral coordinates involve the golden ratio φ = lim(Fₙ₊₁/Fₙ).

\textbf{Status}: STRUCTURAL (mathematical pattern; physical significance unclear)

\begin{center}\rule{0.5\linewidth}{0.5pt}\end{center}

\begin{center}\rule{0.5\linewidth}{0.5pt}\end{center}

\subsection{Appendix F: δ\_CP: The 197° Prediction and the Compactification Correction}\label{appendix-f-ux3b4_cp-the-197-prediction-and-the-compactification-correction}

\subsubsection{F.1 The Original Prediction}\label{f.1-the-original-prediction}

The GIFT prediction δ\_CP = 197° = 7 × 14 + 99 = dim(K₇) × dim(G₂) + H* matches
NuFIT 5.2 (2022: 197° ± 24°) exactly.

NuFIT 6.0 (October 2024) shifted the central value to 177° ± 20°. The 197° prediction
is at the 1σ boundary of this measurement. The NuFIT collaboration notes that δ\_CP
remains one of the least constrained oscillation parameters, with significant sensitivity
to reactor flux model assumptions and the mass ordering.

\textbf{The canonical GIFT prediction remains δ\_CP = 197°.}

\subsubsection{F.2 The Compactification Factor 62/69}\label{f.2-the-compactification-factor-6269}

PSLQ residual analysis (§7.6 of main paper) identifies a potential correction factor:

\[\frac{62}{69} = \frac{\dim(E_8)}{\dim(E_8) + 4\,\dim(K_7)} = \frac{248}{248 + 28} = \frac{248}{276}\]

This factor has a clean structural interpretation: the ratio of gauge degrees of freedom
(dim(E₈) = 248) to total degrees of freedom (248 gauge + 4 × 7 gravitational = 276).

Applied to the original prediction:

\[197 \times \frac{62}{69} = \frac{12214}{69} = 177.01°\]

This matches NuFIT 6.0 to \textbf{0.008\%}. The formula 12214/69 uses only GIFT topological
integers and the compactification ratio.

\subsubsection{F.3 Why We Do Not Adopt It (Yet)}\label{f.3-why-we-do-not-adopt-it-yet}

\begin{enumerate}

\item
  \textbf{Post-hoc identification}: The factor 62/69 was found AFTER the NuFIT 6.0 shift.
  Adopting it would be fitting to the latest central value: the opposite of prediction.
\item
  \textbf{197° remains within 1σ}: The value 197° = 177° + 20° is AT the 1σ boundary.
  One-sigma deviations are expected \textasciitilde32\% of the time. This is not a falsification.
\item
  \textbf{Experimental instability}: δ\_CP shifted 20° between NuFIT 5.2 and 6.0.
  The next update could shift back. DUNE will provide the definitive measurement
  with resolution of a few degrees to \textasciitilde15° (2028--2040).
\item
  \textbf{Formula integrity}: GIFT predictions are derived from topology, not fitted
  to experiment. The 197° formula (3 constants, 2 operations) is cleaner than
  the 177.01° formula (5 constants, 4 operations).
\end{enumerate}

\subsubsection{F.4 Contingency Plan}\label{f.4-contingency-plan}

\textbf{If DUNE confirms δ\_CP ≈ 177° (central value within ±5° of 177°):}
→ Adopt 12214/69 = 177.01° as the primary prediction
→ Interpret 62/69 as a compactification factor (gauge/total DOF ratio)
→ The original 197° becomes the ``bare'' topological value before compactification

\textbf{If DUNE confirms δ\_CP ≈ 197° (central value within ±10° of 197°):}
→ Original prediction vindicated
→ The 62/69 factor was a false lead (PSLQ artifact)

\textbf{If DUNE gives a value inconsistent with both 177° and 197°:}
→ Both formulas fail; the neutrino sector requires revision

\subsubsection{F.5 Summary}\label{f.5-summary}

{\def\LTcaptype{none} % do not increment counter
\begin{longtable}[]{@{}
  >{\raggedright\arraybackslash}p{(\linewidth - 6\tabcolsep) * \real{0.2571}}
  >{\raggedright\arraybackslash}p{(\linewidth - 6\tabcolsep) * \real{0.2000}}
  >{\raggedright\arraybackslash}p{(\linewidth - 6\tabcolsep) * \real{0.2286}}
  >{\raggedright\arraybackslash}p{(\linewidth - 6\tabcolsep) * \real{0.3143}}@{}}
\toprule\noalign{}
\begin{minipage}[b]{\linewidth}\raggedright
Formula
\end{minipage} & \begin{minipage}[b]{\linewidth}\raggedright
Value
\end{minipage} & \begin{minipage}[b]{\linewidth}\raggedright
Status
\end{minipage} & \begin{minipage}[b]{\linewidth}\raggedright
Complexity
\end{minipage} \\
\midrule\noalign{}
\endhead
\bottomrule\noalign{}
\endlastfoot
7×14+99 & 197° & \textbf{PRIMARY} (original prediction) & 3 constants, 2 ops \\
12214/69 & 177.01° & Contingent (pending DUNE) & 5 constants, 4 ops \\
Experimental & 177°±20° & NuFIT 6.0 (1σ band: {[}157°, 197°{]}) & , \\
\end{longtable}
}

\begin{center}\rule{0.5\linewidth}{0.5pt}\end{center}

\subsection{References}\label{references}

\begin{enumerate}

\item
  Joyce, D. D. (2000). \emph{Compact Manifolds with Special Holonomy}. Oxford.
\item
  Atiyah, M. F., Singer, I. M. (1968). \emph{The index of elliptic operators}.
\item
  Particle Data Group (2024). \emph{Review of Particle Physics}. Phys. Rev.~D 110, 030001.
\item
  NuFIT 6.0 (2024). Global neutrino oscillation analysis. www.nu-fit.org.
\item
  Planck Collaboration (2020). Cosmological parameters. A\&A 641, A6.
\item
  T2K, NOvA Collaborations (2025). Nature 638, 534-541.
\end{enumerate}

\begin{center}\rule{0.5\linewidth}{0.5pt}\end{center}

\subsection{Author's Related Works}\label{authors-related-works}

\begin{itemize}
\item
  \textbf{[A]} B. de La Fournière, ``An Explicit Approximate G₂ Metric on a Compact 7-Manifold with Certified Torsion-Free Completion,'' Zenodo \href{https://doi.org/10.5281/zenodo.19892350}{10.5281/zenodo.19892350} (2026).
\item
  \textbf{[B]} B. de La Fournière, ``Spectral Geometry of the G₂-GIFT Manifold: Betti Numbers, KK Spectrum, and Spectral Invariants,'' Zenodo \href{https://doi.org/10.5281/zenodo.19893371}{10.5281/zenodo.19893371} (2026).
\item
  \textbf{[D]} B. de La Fournière, ``An Explicit Closed-Form G₂ Ansatz on a K3-Coassociative Neck with Hyperkähler Rotation and Picard-Lefschetz Wirtinger Certificate,'' Zenodo \href{https://doi.org/10.5281/zenodo.20039066}{10.5281/zenodo.20039066} (2026).
\end{itemize}

\begin{center}\rule{0.5\linewidth}{0.5pt}\end{center}

\emph{GIFT Framework: Supplement S2}
\emph{Complete Derivations: 33 Type I Relations}


\section*{References}
\addcontentsline{toc}{section}{References}

\input{gift_3.4_bib}

\noindent\rule{\textwidth}{0.2pt}

\section*{Author's Related Works}
\begin{itemize}
\item \textbf{[A]} B.~de La Fourni\`ere, ``An Explicit Approximate $\mathrm{G}_2$ Metric on a Compact 7-Manifold with Certified Torsion-Free Completion,'' Zenodo \href{https://doi.org/10.5281/zenodo.19892350}{10.5281/zenodo.19892350} (2026).
\item \textbf{[B]} B.~de La Fourni\`ere, ``Spectral Geometry of the $\mathrm{G}_2$-GIFT Manifold: Betti Numbers, KK Spectrum, and Spectral Invariants,'' Zenodo \href{https://doi.org/10.5281/zenodo.19893371}{10.5281/zenodo.19893371} (2026).
\item \textbf{[C]} B.~de La Fourni\`ere, ``Newton-Kantorovich Certificate for the K3 Donaldson Embedding in the $\mathrm{G}_2$-GIFT Metric,'' Zenodo \href{https://doi.org/10.5281/zenodo.19708916}{10.5281/zenodo.19708916} (2026).
\item \textbf{[D]} B.~de La Fourni\`ere, ``An Explicit Closed-Form $\mathrm{G}_2$ Ansatz on a K3-Coassociative Neck with Hyperk\"ahler Rotation and Picard-Lefschetz Wirtinger Certificate,'' Zenodo \href{https://doi.org/10.5281/zenodo.20039066}{10.5281/zenodo.20039066} (2026).
\item \textbf{[essay]} B.~de La Fourni\`ere, \emph{Orientation, not ontology} (companion essay), \url{https://giftheory.substack.com/p/orientation-not-ontology} (2026).
\end{itemize}

\noindent\rule{\textwidth}{0.2pt}

\begin{center}
\emph{GIFT Framework --- Supplement S2 --- Complete Derivations: 33 Type I Relations}
\end{center}

\end{document}
